\documentclass[10pt,a4paper]{article}

% ----------------- Paquetes -------------------%
\usepackage[T1]{fontenc}
\usepackage[utf8]{inputenc}
\usepackage[a4paper,margin=2.5cm]{geometry}
\usepackage{mathptmx}             
\usepackage{changepage} 
\usepackage{setspace}
\usepackage[spanish]{babel}
\usepackage[labelfont=bf,labelsep=space]{caption}
\usepackage{amssymb}
\usepackage{amsmath}
\usepackage{ebgaramond}
\usepackage{placeins}
\usepackage{etoolbox}
\usepackage{setspace}
\usepackage{parskip} 
\usepackage{tcolorbox}
\usepackage{needspace}
\usepackage{xparse}
\usepackage{ocgx2}
\usepackage{hyperref}
\usepackage{hypcap}
\usepackage{multirow}
\usepackage{soul}\sethlcolor{yellow}% resaltado amarillo: \hl{texto} (Panel TCEE, ⌃H)



%%%%%%%%%%%%%%%%%%%%%%%%%%%%%%%%%%%%%%%%%%%%%%%%%%%%%%%%%%%%%%%%
% --- Fija primero tus valores globales "buenos" ---
\setstretch{1.2}                 % Interlineado global
\setlength{\parskip}{0.7\baselineskip}
\setlength{\parindent}{0pt}
\newlength{\GlobalParskip}
\setlength{\GlobalParskip}{\parskip}

\newcounter{eqnum}[subsection]
\renewcommand{\theeqnum}{\arabic{eqnum}}
\newcounter{alphaitem}
\newcounter{numitem}

%% ------------------- Recuadros----------------------%%
\newcommand{\puntosclave}[1]{%
  \begin{tcolorbox}[
    colframe=black,
    title={Puntos clave},
    before upper={\setlength{\parindent}{0.5cm}\setlength{\parskip}{0.5\baselineskip}}
  ]
  #1
  \end{tcolorbox}
}

\newcommand{\modificaciones}[1]{
  \begin{tcolorbox}[
    colback=white,
    colframe=red,
    title={Modificaciones a realizar},
    before upper={\setlength{\parindent}{0.5cm}\setlength{\parskip}{0.5\baselineskip}}
  ]
  #1
  \end{tcolorbox}
}

%% ------------------- Preguntas TEST----------------------%%
% Opciones: radio (single-choice)
\NewDocumentCommand{\OptRadio}{m m m}{%
  \noindent\CheckBox[radio,name=#1,value=#2,width=1.2ex,height=1.2ex]{}%
  \;\textbf{#2.}~#3\par
}

% Opciones: checkbox (multi-choice)
\NewDocumentCommand{\OptCheck}{m m}{%
  \noindent\CheckBox[name=#1,width=1.2ex,height=1.2ex,bordercolor={0 0 0}]{}%
  \;#2\par
}

% Pregunta single-choice (radio)
% Args: 1=id, 2=enunciado, 3=opciones, 4=solución+justificación, 5=imagen (o {}), 6=origen
\NewDocumentCommand{\PreguntaTestSingle}{m m m m m m}{%
  \par\medskip
  \noindent\textbf{#1.}~#2\par\smallskip

    % Imagen (si no es {})
    \IfBlankTF{#5}{}{%
      \begin{center}
        \includegraphics[width=0.75\linewidth]{#5}
      \end{center}
    }

    \begin{Form}
      #3
    \end{Form}

    \par\smallskip
    \switchocg{sol-#1}{\fbox{Mostrar / ocultar solución}}%
    \hfill{\footnotesize #6}\par\smallskip

    \begin{ocg}{Solución}{sol-#1}{0}
      \noindent #4
    \end{ocg}
  \par\medskip
}

% Pregunta multi-choice (checkboxes)
\NewDocumentCommand{\PreguntaTestMulti}{m m m m m m}{%
  \par\medskip
  \noindent\textbf{#1.}~#2\par\smallskip

    % Imagen (si no es {})
    \IfBlankTF{#5}{}{%
      \begin{center}
        \includegraphics[width=0.75\linewidth]{#5}
      \end{center}
    }
    \begin{Form}
      #3
    \end{Form}

    \par\smallskip
    \switchocg{sol-#1}{\fbox{Mostrar / ocultar solución}}%
    \hfill{\footnotesize #6}\par\smallskip

    \begin{ocg}{Solución}{sol-#1}{0}
      \noindent #4
    \end{ocg}
  \par\medskip
}

%% ------------------- CITA ----------------------%%
% \begin{cita}[Autor][Año][Obra] texto \end{cita}: cita sangrada y, debajo a la derecha, «AUTOR (Año), Obra». Los tres datos son opcionales.
% En el Panel TCEE: escribir \cita e Intro (el cursor va al texto; Tab pasa a Autor, Año y Obra).
\NewDocumentEnvironment{cita}{ O{} O{} O{} +b }{%
  \begin{adjustwidth}{2cm}{0pt}
  \raggedright
  #4\par
  \raggedleft
  \if\relax\detokenize{#1#2#3}\relax
    % sin firma
  \else
    #1%
    \if\relax\detokenize{#2}\relax\else\if\relax\detokenize{#1}\relax\else\space\fi(#2)\fi
    \if\relax\detokenize{#3}\relax\else\if\relax\detokenize{#1#2}\relax\else, \fi\textit{#3}\fi
  \fi
  \end{adjustwidth}
}{}



%% ------------------- Listas----------------------%%
    % --- Lista alfabética ---
    \newenvironment{la}
      {\begin{list}{\alph{alphaitem})}{%
          \usecounter{alphaitem}%
          \setlength{\leftmargin}{1cm}%
          \setlength{\parskip}{3pt}% 
          \setlength{\itemsep}{0pt}% ← espacio entre ítems
          \setlength{\parsep}{0pt}%  ← párrafos dentro de un ítem
      }}
      {\end{list}}
      
    % --- Lista numérica ---
    \newenvironment{lnum}
      {\begin{list}{\arabic{numitem}.}{%
          \usecounter{numitem}%
          \setlength{\leftmargin}{1cm}%
          \setlength{\parskip}{3pt}% 
          \setlength{\itemsep}{0pt}% ← espacio entre ítems
          \setlength{\parsep}{0pt}%  ← párrafos dentro de un ítem
      }}
      {\end{list}}

%% ------------------- Preguntas Test ----------------------%%
      \usepackage{xparse} % si ya lo tienes, no lo repitas

    \NewDocumentEnvironment{test}{m m o}{%
      \begin{tcolorbox}[
        colback=white,
        colframe=black!15,
        boxrule=0.6pt,
        arc=2mm,
        left=2mm,right=2mm,top=2mm,bottom=2mm
      ]
      \centering
      \includegraphics[width=\linewidth]{#1}
      \end{tcolorbox}
    
      \vspace{2mm}
    
      % Caja SOLO para texto (SÍ partible y mantiene márgenes al partir)
      \begin{tcolorbox}[
        enhanced,
        breakable,
        colback=black!3,
        colframe=black!25,
        boxrule=0.6pt,
        arc=2mm,
        left=2mm,right=2mm,top=1.5mm,bottom=1.5mm
      ]
      \textbf{Respuesta:} \textbf{#2}%
      \IfValueT{#3}{\par\smallskip \textbf{Justificación:} #3}
      \end{tcolorbox}
    }{}
      
% ------------------- Imágenes----------------------%
    \graphicspath{{Img/}}% Rutas por defecto 
    % --- Imagen adaptativa ---
    % Registros de dimensión definidos una sola vez
    \newdimen\imgcap
    \newdimen\imgmin
    \newdimen\imgmax
    \newdimen\imgremain
    \newdimen\imgh
    \newdimen\imgneed
    
    \newcommand{\imagenfit}[3]{%
      \addvspace{10pt}% separación antes
      \par\begingroup
        % Parámetros de composición (asignaciones, no \newdimen)
        \imgcap=1\baselineskip            % alto estimado de título+fuente
        \imgmin=0.15\textheight
        \imgmax=0.15\textheight
        \imgremain=\pagegoal
        \ifdim\pagegoal=\maxdimen \imgremain=0pt\fi
        \advance\imgremain by -\pagetotal
        \advance\imgremain by -\imgcap
        % Decisión de altura destino
        \ifdim\imgremain<\imgmin
          \newpage
          \imgh=0.20\textheight
        \else
          \imgh=\imgremain
          \ifdim\imgh>\imgmax \imgh=\imgmax \fi
        \fi
        % Reservar espacio para evitar cortes del bloque completo
        \imgneed=\imgh \advance\imgneed by \imgcap
        \Needspace{\imgneed}
        % Cuerpo
        \noindent\begin{minipage}{\textwidth}
          \centering
          \captionof{figure}{\textemdash\ #2}\par\vspace{0.3em}% título arriba
          \includegraphics[width=\linewidth,height=\imgh,keepaspectratio]{\detokenize{#1}}\par
          \vspace{0.3em}\small\textit{Fuente: }#3% fuente abajo
        \end{minipage}%
      \endgroup
      \par\addvspace{10pt}%
    }

%------------ Bloque de RECUADRO ESPECIAL -------------%
    \newenvironment{specialblock}[1]{%
      \begin{quote} % crea sangría a izquierda y derecha
      \small % texto más pequeño
      \noindent\textbf{#1}\\[-0.3em] % título en negrita
      \rule{\linewidth}{0.4pt}\\[0.6em]}{
      \end{quote}}
      
%-------------------------- Portada -----------------------%
\newcommand{\oldtitle}[1]{\def\OldTitle{#1}}
\newcommand{\changerationale}[1]{\def\ChangeWhy{#1}}

\newcommand{\changebox}{%
\begin{tcolorbox}[title={Historial de cambios del tema}]
\textbf{Título anterior:} \OldTitle\par
\textbf{Motivo del cambio:} \ChangeWhy
\end{tcolorbox}
}

%-------------------------- Author Font -----------------------%
    \newcommand{\authorfont}[1]{{\fontfamily{EBGaramond-TLF}\selectfont #1}}

%-------------------------- Anexos -----------------------%
    \makeatletter
    \newcounter{anexo}
    \renewcommand{\theanexo}{\Roman{anexo}}
    \newcommand{\anexo}[1]{%
      \clearpage
      \phantomsection
      \refstepcounter{anexo}%
      \noindent\textbf{Anexo \theanexo.- }#1\par\medskip
      \addcontentsline{stc}{section}{Anexo \theanexo.- #1}%
    }
    \makeatother

    %-------------------------- Lista de figuras (personal) -----------------------%
\makeatletter
\newcommand{\MyListOfFigures}{%
  \clearpage
  \phantomsection
  {\centering\bfseries\Large Índice de figuras\par}%
  \medskip
  % Entrada en nuestro índice de contenido (stc)
  \addcontentsline{stc}{section}{Índice de figuras}%
  % Recuperación del listado (sin cabecera propia)
  \begingroup
    \parindent=0pt\parskip=2pt
    \@starttoc{lof}%
  \endgroup
}
\makeatother

%-------------------------- Bibliografía -----------------------%
\makeatletter
% Contador para las referencias
\newcounter{myref}
% Cabecera de la bibliografía, entra en el índice 'stc'
\newcommand{\MyBibliography}{%
  \clearpage
  \phantomsection
  {\centering\bfseries\Large Bibliografía y referencias\par}%
  \medskip
  \addcontentsline{stc}{section}{Bibliografía y referencias}%
  \setcounter{myref}{0}%
}
\newcommand{\MyBibItem}[4]{%
  \refstepcounter{myref}%
  \noindent[\themyref]\ #1.\ \emph{#2}. #3. #4\par
}
\makeatother

%--------- Establecer cuatro niveles de índice --------%
    \setcounter{tocdepth}{4}   
    \setcounter{secnumdepth}{4}% numera hasta \paragraph

%%%%%%%%%%%%%%%%%%%%%%%%%%%%%%%%

\makeatletter

    %--------- Interlineado Global --------%
    \edef\GlobalBaseLineStretch{\baselinestretch}
    
    %--------- Espaciado de seccinones --------%
    \renewcommand\section{\@startsection{section}
      {1}{\z@}
      {2.5ex \@plus 1ex \@minus .2ex} % espacio antes
      {1.5ex \@plus .2ex}             % espacio después
      {\normalfont\Large\bfseries}    % formato del título
    }
    \renewcommand\subsection{\@startsection{subsection}{2}{\z@}
      {3ex \@plus 0.8ex \@minus .2ex}
      {1.5ex \@plus .2ex}
      {\normalfont\large\bfseries}
    }
    \renewcommand\subsubsection{\@startsection{subsubsection}{3}{\z@}
      {3ex \@plus 0.5ex \@minus .2ex}
      {1ex \@plus .2ex}
      {\normalfont\normalsize\bfseries}
    }
    \renewcommand\paragraph{\@startsection{paragraph}{4}{\z@}%
    {-2ex \@plus -0.5ex \@minus -.2ex}% espacio antes
    {0.8ex \@plus .2ex}% espacio después
    {\normalfont\normalsize\itshape}}% ← cursiva en lugar de negrita

    %--------- Ecuaciones --------%   
    % Variables globales para almacenar títulos y notas
    \gdef\leq@current@section{}%
    \gdef\leq@current@subsection{}%
    \gdef\leq@last@printed@section{}%
    \gdef\leq@last@printed@subsection{}%
    
    % Comando para añadir nota (invisible en el cuerpo)
    \newcommand{\eqnote}[1]{%
      \addcontentsline{leq}{leqnote}{#1}%
    }
    
    % Comando para ecuaciones
    \newcommand{\eqblock}[2]{%
      % Verificar si necesitamos imprimir el título de sección
      \ifx\leq@current@section\leq@last@printed@section\else
        \addcontentsline{leq}{leqsection}{\leq@current@section}%
        \global\let\leq@last@printed@section\leq@current@section
        \gdef\leq@last@printed@subsection{}% Reset subsección al cambiar sección
      \fi
      % Verificar si necesitamos imprimir el título de subsección
      \ifx\leq@current@subsection\leq@last@printed@subsection\else
        \ifx\leq@current@subsection\@empty\else
          \addcontentsline{leq}{leqsubsection}{\leq@current@subsection}%
          \global\let\leq@last@printed@subsection\leq@current@subsection
        \fi
      \fi
      % Ecuación
      \refstepcounter{eqnum}%
      \begin{equation*}
        #1\tag{E.\theeqnum}
      \end{equation*}%
      \addcontentsline{leq}{leqt}{[E.\theeqnum] -- #2}%
      \addcontentsline{leq}{leqm}{#1}%
    }
    
    %%% ----- Índice de ecuaciones ----------%%%
    % Tamaño para []
    \newcommand{\leqIndexFont}{\normalsize}
    \newcommand{\leqTitleFont}{\tiny}
    \newcommand{\leqNoteFont}{\tiny}
    % Cabecera de sección 
    \newcommand*\l@leqsection[2]{%
      \addvspace{25pt}% Mayor separación antes de nueva sección
      \noindent\leqIndexFont
      \makebox[\linewidth]{\rule{.38\linewidth}{0.4pt}\ \ \textbf{  \MakeUppercase{#1}}\ \ \rule{.38\linewidth}{0.4pt}}%
      \par\vspace{-10pt}
    }
    % Cabecera de subsección 
    \newcommand*\l@leqsubsection[2]{%
      \par\addvspace{15pt}%
      \noindent\leqIndexFont #1
      \par\vspace{-7pt}
    }
    % Título de ecuación 
    \newcommand*\l@leqt[2]{%
    \par\addvspace{10pt}%
      \par\noindent\leqTitleFont #1\par
    }
    % Miniatura de ecuación
    \newcommand*\l@leqm[2]{%
      \vspace{0.3\baselineskip}%
      \par\scriptsize\noindent\hspace*{1.5em}\(#1\)\par%
    }
    % Nota de ecuación 
    \newcommand*\l@leqnote[2]{%
      \vspace{0.2\baselineskip}%
      \par\noindent\leqNoteFont\hspace*{1.5em}\textbf{Nota.--} #1\par\vspace{0.5\baselineskip}%
    }
    % Generación del índice
    \newcommand{\mylistofequations}{%
      \clearpage
      \phantomsection
      % Título visual de la lista (ajústalo a tu gusto):
      {\centering\bfseries\Large Índice de ecuaciones\par}
      % Entrada en nuestro índice 'stc':
      \addcontentsline{stc}{section}{Índice de ecuaciones}%
      % Llamada a tu lista real (usa el nombre exacto de tu macro):
      \@starttoc{leq}%
    }
    
    %--------- Captura de títulos de sección --------%
    \let\leq@original@section\section
    \let\leq@original@subsection\subsection
    % Flag para saber si estamos en una sección asterisco
    \newif\ifLeqStarSection
    \LeqStarSectionfalse
    
    % Redefinir \section CON CONTROL DE ASTERISCO
    \renewcommand{\section}{%
      \@ifstar{\leq@section@star}{\@dblarg\leq@section@nostar}%
    }
    \def\leq@section@star#1{%
      \global\LeqStarSectiontrue
      \gdef\leq@current@section{#1}%
      \gdef\leq@current@subsection{}%
      \leq@original@section*{#1}%
      \global\LeqStarSectionfalse
    }
    \def\leq@section@nostar[#1]#2{%
      \global\LeqStarSectionfalse
      \gdef\leq@current@section{#2}%
      \gdef\leq@current@subsection{}%
      \leq@original@section[#1]{#2}%
    }
    % Redefinir \subsection
    \renewcommand{\subsection}{%
      \@ifstar{\leq@subsection@star}{\@dblarg\leq@subsection@nostar}%
    }
    \def\leq@subsection@star#1{%
      \global\LeqStarSectiontrue
      \gdef\leq@current@subsection{#1}%
      \leq@original@subsection*{#1}%
      \global\LeqStarSectionfalse
    }
    \def\leq@subsection@nostar[#1]#2{%
      \global\LeqStarSectionfalse
      \gdef\leq@current@subsection{#2}%
      \leq@original@subsection[#1]{#2}%
    }

    % ----- Restauración de valores Globales de estilo ----------%
    % Flags
    \newif\ifInFootnote
    \InFootnotefalse
    \pretocmd{\@footnotetext}{\InFootnotetrue}{}{}
    \apptocmd{\@footnotetext}{\InFootnotefalse}{}{}
    % Profundidad de listas (soporta anidadas)
    \newcount\MyListDepth \MyListDepth=0
    \AtBeginEnvironment{la}{\advance\MyListDepth by 1}
    \AfterEndEnvironment{la}{\advance\MyListDepth by -1}
    \AtBeginEnvironment{lnum}{\advance\MyListDepth by 1}
    \AfterEndEnvironment{lnum}{\advance\MyListDepth by -1}
    \AtBeginEnvironment{itemize}{\advance\MyListDepth by 1}
    \AfterEndEnvironment{itemize}{\advance\MyListDepth by -1}
    \AtBeginEnvironment{enumerate}{\advance\MyListDepth by 1}
    \AfterEndEnvironment{enumerate}{\advance\MyListDepth by -1}
    % Restauración diferida: hazlo al INICIO del próximo párrafo real
    \newif\if@pendingRestore
    \@pendingRestorefalse
    \newcommand{\RequestRestore}{\global\@pendingRestoretrue}
    \everypar{%
      \if@pendingRestore
        \global\@pendingRestorefalse
        \ifInFootnote\else
          \ifnum\MyListDepth=0
            \setlength{\parskip}{\GlobalParskip}%
            \setstretch{\GlobalBaseLineStretch}%
          \fi
        \fi
      \fi
    }
    % Al final de bloques:
    \AfterEndEnvironment{equation}{\RequestRestore}
    \AfterEndEnvironment{equation*}{\RequestRestore}
    \AfterEndEnvironment{la}{\RequestRestore}
    \AfterEndEnvironment{lnum}{\RequestRestore}
    % Al final de macros:
    \apptocmd{\eqblock}{\RequestRestore}{}{}
    \apptocmd{\imagenfit}{\RequestRestore}{}{}
    
\makeatother

% ===================== ToC propio con 4 niveles (único bloque) =====================
\makeatletter

% --- Escritores: añadimos una línea al .stc en cada cabecera numerada ---
% Guardamos las versiones actuales para llamar al comportamiento normal
\let\MyTOC@orig@section\section
\let\MyTOC@orig@subsection\subsection
\let\MyTOC@orig@subsubsection\subsubsection
\let\MyTOC@orig@paragraph\paragraph

% SECTION
\renewcommand{\section}{%
  \@ifstar{\MyTOC@section@star}{\@dblarg\MyTOC@section@nostar}%
}
\def\MyTOC@section@star#1{%
  \MyTOC@orig@section*{#1}% sin entrada al .stc
}
\def\MyTOC@section@nostar[#1]#2{%
  \MyTOC@orig@section[{#1}]{#2}% comportamiento normal (escribe al .toc)
  % Escribimos también nuestra línea al .stc (texto con \numberline)
  \addcontentsline{stc}{section}{\protect\numberline{\thesection}#2}%
}

% SUBSECTION
\renewcommand{\subsection}{%
  \@ifstar{\MyTOC@subsection@star}{\@dblarg\MyTOC@subsection@nostar}%
}
\def\MyTOC@subsection@star#1{%
  \MyTOC@orig@subsection*{#1}%
}
\def\MyTOC@subsection@nostar[#1]#2{%
  \MyTOC@orig@subsection[{#1}]{#2}%
  \addcontentsline{stc}{subsection}{\protect\numberline{\thesubsection}#2}%
}

% SUBSUBSECTION
\renewcommand{\subsubsection}{%
  \@ifstar{\MyTOC@subsubsection@star}{\@dblarg\MyTOC@subsubsection@nostar}%
}
\def\MyTOC@subsubsection@star#1{%
  \MyTOC@orig@subsubsection*{#1}%
}
\def\MyTOC@subsubsection@nostar[#1]#2{%
  \MyTOC@orig@subsubsection[{#1}]{#2}%
  \addcontentsline{stc}{subsubsection}{\protect\numberline{\thesubsubsection}#2}%
}

% PARAGRAPH
\renewcommand{\paragraph}{%
  \@ifstar{\MyTOC@paragraph@star}{\@dblarg\MyTOC@paragraph@nostar}%
}
\def\MyTOC@paragraph@star#1{%
  \MyTOC@orig@paragraph*{#1}%
}
\def\MyTOC@paragraph@nostar[#1]#2{%
  \MyTOC@orig@paragraph[{#1}]{#2}%
  % Si no numeras los \paragraph, cambia \theparagraph por texto plano sin \numberline
  \addcontentsline{stc}{paragraph}{\protect\numberline{\theparagraph}#2}%
}

% --- Impresor del índice propio (lee el .stc) ---
\newcommand{\MyTOC}{%
  \par\bigskip
  {\centering\bfseries\Large Índice de contenido\par}%
  \medskip
  \begingroup
    \parindent=0pt\parskip=2pt
    \@starttoc{stc}%
  \endgroup
}

\makeatother
% ===================== fin ToC propio =====================


%%%%%%%%%%%%%%


% --- Datos del tema ---
\title{Tema 3.A.15 -- La empresa: el tamaño eficiente y sus límites.\\
\large Mención especial de la economía de los costes de transacción. La Teoría de la Organización Industrial: barreras a la entrada y mercados impugnables.}
\author{Marcelo Ortega}
\date{Diciembre 2025}
\newcommand{\TemaCode}{3A15}

\oldtitle{Tema 3.A.11. Análisis teórico de la empresa en el marco de la teoría de la organización industrial}
\changerationale{
    Se desarrolla y aclara el título para que el opositor evite una narrativa histórica de las teorías y se centre en los desarrollos de mayor relevancia, en concreto, se podrían incluir aquellos por los que Tirole y Hart y Holmstrom recibieron los Premios Nobel en 2014 y 2016 respectivamente.
}

\begin{document}


\maketitle
\changebox

\newpage

% --- Página de resumen y modificaciones ---
\puntosclave{ %% Escribir puntos clave Apuntes CECO
    Se debe hacer hincapié tanto en la economía de los costes de transacción como en la teoría de los mercados contestables de Baumol. A pesar de que el título haga una referencia explícita a las "barreras de entrada", dicha referencia debería entenderse en el contexto de explicar la teoría de BAUMOL (y, si se quiere, nuevas aplicaciones en los mercados digitales, como se hace en este tema). Así, no deberían incluirse en este tema modelos de barreras de entrada que típicamente se incluyen a la hora de explicar la Teoría del Oligopolio en el largo plazo.
    
    La estructura que debe tener el tema, explicando por un lado la teoría de la empresa y, por otro lado, la teoría de los mercados o Teoría de la Organización Industrial.
    }
\vspace{1cm}

\modificaciones{

    \textcolor{red}{\textbf{OJO}} Revisar y comparar TODA la parte de la NTOI con la de Podr de Mercado y Regulación, se trata en realidad del mismo contenido. Primer apartado tema 20 y parte de NTOI de este tema son iguales.

    \textbf{NOTA (04/10/2026):} Revisar (Rosell): https://competition-policy.ec.europa.eu/document/download/dc48d9d4-8151-454c-9bfc-c69af2c96e86\_en?filename=kd0126016enn\_Impact\_mergers\_innovation\_and\_markups.pdf
}

\newpage
\MyTOC
\newpage

\mylistofequations
\clearpage

\MyListOfFigures
\clearpage


\pagenumbering{arabic}
\setcounter{page}{1}


\section{Introducción}\label{intro}


\subsection{Contextualización}\label{context}


\subsubsection{Relevancia}\label{importance}


\subsubsection{Historia}\label{history}



\subsubsection{Problemática}\label{problem} 




\subsection{Estructura de la exposición}\label{structure}







\clearpage
\section{Teorías sobre la existencia y estructura interna de la Empresa}
\puntosclave{
    La economía neoclásica concibe a la empresa como una tecnología que convierte factores de producció n en bienes o servicios, sin más condicionantes ni análisis. Los conceptos de economías de escala y economías de alcance explican tanto el tamaño de la empresa como la configuración ó ptima de la estructura de mercado. No obstante, no es capaz de explicar otros aspectos de la realidad empresa rial, como los cont ratos de suministro.
    
	La Teoría de Contratos, basada en la idea de Coase de los costes de transacción, pe rmite añadir muchos matices sobre la concepción y el comportamiento de las empresas. En función de que los cont ratos puedan prever cualquier tipo de contingencia se diferencia entre las teorías de contratos completos e incompletos. Entre otros comportamientos, dichas teorías permiten explicar el dilema to make vs. to buy.
    
	Las plataformas digitales son un nuevo tipo de empresa, con unas características propias (efectos de red, economías de escala y costes de cambio) que les permite separarse de la regla neoclásica de fijación de precios en función de las elasticidades inversas.
}
	
Desde que existe la producción en una economía, los agentes econ ómicos interaccionan entre sí para intercambiar diferentes bienes. Para llevar a cabo estas transacciones de forma fluida se necesitan ciertos instrumentos e instituciones que permitan coordinarlas. Existen dos instituciones de coordinación fundamentales: el mercado y las organizaciones (por ejemplo, las empresas). El mercado utiliza el precio, mientras que las organizaciones utilizan su influencia, entre otros elementos, para realizar la coordinación entre los agentes. Desde una perspectiva complementaria, puede decirse que los intercambios en una economía pueden ser "descentralizados", cuando se realizan en el mercado (esto es, entre distintos compradores y vendedores) o centralizados, cuando se realizan en el seno de una organización gobernada por una serie de reglas (por ejemplo, la prestación de un servicio de un trabajador a su directivo a cambio de un salario).

\subsection{Modelo baseline: Enfoque tecnológico}
El enfoque tecnológico es el propio de la teoría neoclásica tradicional, que ya habíamos avanzado: la empresa como una caja negra que combina factores productivos para producir bienes y servicios, es decir, la empresa como relación puramente técnica entre inputs y outputs.
El enfoque tecnológico ofrecerá explicaciones a cuestiones relacionadas con el tamaño de la empresa, la configuración de la industria, o el grado de integración.

\subsubsection{Tamaño}
El tamaño en el enfoque tecnológico vendrá determinado por la estructura de costes que resulte de su tecnología. Así, una empresa crecerá hasta agotar sus economías de escala, es decir, la empresa alcanzará el tamaño de planta que minimice sus costes medios.

\subsubsection{Configuración de la Industria}
La estructura de costes de las empresas puede predeterminar la configuración de la industria en la que se mueven. Así, relacionado con las economías de escala que mencionábamos, podemos analizar la subaditividad de costes, que es aquella estructura de costes para la cual resulta menos costoso concentrar toda la producción en una sola empresa que distribuirla entre varias, por lo que la estructura eficiente de la industria no será la de competencia perfecta.

$$C\left(\sum_{i=1}^nq_i\right)<\sum_{i=1}^nC_i(q_i)$$
 
\subsubsection{Grado de integración}
El enfoque tecnológico también ofrece una explicación al grado de integración de las empresas: éstas abarcarán más o menos actividades en función de la intensidad de las economías de alcance, de forma que las empresas estarán muy integradas si la producción conjunta de varios bienes es menos costosa que su producción por distintas empresas.\footnote{%
    Entre los factores que explican las economías de alcance están los inputs públicos, que una vez utilizados para un producto, sirven para producir otro sin coste alguno y la existencia de inputs aplicables para producir distintos productos como el capital. Un abogado que presta asesoría fiscal y defensa en juicio o la industria alimentaria.
} [Se puede hablar del concepto de Relación Marginal de Transformación o Curva de Transformación]	

\href{https://policonomics.com/es/curva-transformacion/}{Curva de Transformación}

En definitiva, La tecnología y la demanda nos permiten explicar el tamaño de la empresa (el que permita agotar íntegramente las economías de escala y de alcance) y la configuración eficiente de la industria, que será aquélla que permita sostener un equilibrio de mercado a los menores costes posibles.\footnote{%
    Ahora bien, si consideramos la existencia de tecnologías alternativas (como un exceso de capacidad en la función de costes, teoría de los costes de STIGLER), se puede llegar a conclusiones distintas a la anterior pudiendo existir una horquilla de dimensiones óptimas entre un mínimo y un máximo.
} 

\subsection{Modelo contractual}
No obstante, el enfoque tecnológico es insuficiente. Entre otros no llega a explicar: (i) por qué la solución más conveniente es explotar economías de escala dentro de la misma empresa en vez de suscribir un contrato entre empresas legalmente separadas;\footnote{%
    A pesar de sus limitaciones, la visión tecnológica de la empresa es esencial, porque permite analizar un aspecto crucial de la actividad productiva: la empresa como mecanismo de maximización de beneficios y/o minimización de costes [\authorfont{Ver Tema 3.A.12}] (ii)
} lo que ocurría dentro de la caja negra, que era ignorado; y, (iii) omite todos los aspectos relacionados con cuestiones organizativas internas.\footnote{%
    TIROLE califica este enfoque como la empresa como sinergia estática siendo un enfoque reduccionista (no tiene en cuenta otras cuestiones y actividades de la empresa como la iniciativa empresarial en el know how, marketing empresarial o actividad gerencial).
}

El economista estadounidense RONALD COASE, fue el primero en criticar la escasa fundamentación de la teoría de la empresa neoclásica en su célebre artículo \textit{The Nature of the Firm} (1937), basándose en las constataciones que FRANK KNIGHT había realizado unos años antes.\footnote{%
    La teoría de KNIGHT expuesta en \textit{Risk, Uncertainty and Profit} (1921) supone una crítica a la teoría neoclásica, al partir de un mundo incierto, dinámico y cambiante; por lo que se suele categorizar más bien como una teoría del beneficio o del empresario. Según KNIGHT, en condiciones de competencia perfecta no habría beneficio, porque el residuo obtenido por el empresario sólo podría producirse si la demanda y la oferta son desiguales. Es decir, el beneficio se origina en la incertidumbre. Tanto el precio de las materias primas como el de los costes de producción suponen la principal fuente de incertidumbre. Éste autor pensaba que toda empresa es incierta, y el empresario -que ejerce el control final- tiene que asumir las consecuencias y reducir los efectos de esta incertidumbre. Así, el dirigente de la empresa es responsable de la misma y debe desarrollar los juicios adecuados para hacer frente a la incertidumbre. PIGOU añadiría después que el beneficio puede considerarse como un ingreso que se debe al empresario por tomar sus decisiones en tal incertidumbre.
}

En síntesis, COASE se pregunta por qué existen las empresas (o entornos organizativos) si, como se ha expuesto anteriormente, los agentes económicos pueden elegir entre transaccionar en el mercado o dentro de una organización o empresa. Al fin y al cabo, la actividad productiva podría organizarse en torno a transacciones puntuales materializadas en contratos entre los diferente agentes, y la empresa, como organización, no tendría por qué tener cabida en este entorno. 

\subsubsection{Costes: COASE (1937)}
La razón principal por la que puede ser más eficiente establecer una empresa es porque el uso del mecanismo de precios (esto es, de contratar en el mercado) tiene costes, los cuales pueden ser más elevados que si una organización (empresa) los coordinara por sí sola.\footnote{%
    Sin la empresa, el más simple de los productos llevaría consigo una compleja red de contratos entre productores aislados. Por ejemplo, si un bien conlleva 5 procesos distintos, en ausencia de una empresa, cada productor individual tendría que firmar un contrato con los otros 4, lo que supone la firma de 10 contratos. Por el contrario, si uno de los productores se constituye como empresa y emplea al resto de productores, sólo se requerirán 4 contratos.
} 

Para COASE, estos costes incluían los costes de encontrar los precios adecuados en el mercado, los costes de la negociación y la firma del contrato. Por lo tanto, la empresa comparará los costes de realizar esas actividades dentro de la propia empresa con los derivados de realizar esa actividad en el mercado, de manera que la empresa crecerá hasta que los costes de coordinación internos (relacionadas, por lo general, con las deseconomías gerenciales) sean iguales a los costes de transacción: ${CMg}_T^{\text{emp}} = {CMg}_T^{\text{mdo}}$. La jerarquía empresarial reduce los costes de transacción y la incertidumbre.\footnote{%
    En una empresa, la asignación se determina a través de una estructura de decisión jerarquizada (las “islas de planificación consciente” en un océano de cooperación inconsiciente), en las que el mecanismo de precios está suspendido (mecanismo propio de la teoría neoclásica en la cual la interacción de consumidores y empresas determina los precios de mercado).
} 

No obstante, debemos acudir a APPEL (1996)\footnote{%
    Retoma el enfoque de WILLIAMSON, quien distinguía entre ex ante y ex post, pero lo sistematiza en una estructura más precisa de categorías.
\begin{la}
    \item Ex ante: los que se producen antes de que exista o se ejecute el contrato; corresponden al proceso de formación y estructuración del acuerdo; y,
    \item Ex post: los que aparecen después de celebrado el contrato, durante su ejecución, supervisión o eventual renegociación.
\end{la}
} para tener una mayor sistematización de los costes, quien amplía el análisis de los transaction costs mostrando que no son un bloque homogéneo, sino que pueden descomponerse según el momento temporal (ex ante o ex post) en el que se generan y las funciones que cumplen dentro del proceso contractual, más allá del precio de la misma.

\paragraph{Costes de búsqueda e información}
Para el autor, estos costes preceden a la transacción efectiva y se orientan a la reducción de la asimetría informativa inicial.\footnote{%
    Más allá de la relevancia microeconómica de estos costes en la teoría de la empresa, recuérdese la importancia de los mismos a la hora de analizar el ciclo macroeconómico de un país y las posibles fuentes de fricciones (Modelos de búsqueda en el mercado laboral, teoría de la demanda del dinero). Lo mismo podría decirse de los costes de ajuste (Demanda de trabajo intertemporal, Demanda de inversión) o de información (Expectativas semirracionales).
} Implican la identificación de contrapartes, la recopilación de información sobre precios, calidad, reputación o condiciones del mercado. Su función es reducir la incertidumbre inicial y asegurar una elección racional del socio o del activo.

\paragraph{Costes de decisión y negociación}
El autor los describe como costes de estructuración de la transacción, que son también ex ante porque se incurren antes de que el contrato sea operativo, buscando prevenir conflictos futuros previos a su ejecución. Incluyen el tiempo, los recursos y la asistencia legal o técnica necesarios para negociar los términos del contrato, redactarlo y acordar salvaguardas.

\paragraph{Costes de control (incluida la desincentivación), adaptación (o ajuste) y cumplimiento (o ejecución)}
El autor los asocia a la “dinámica contractual”, enfatizando que incluso con un contrato cuidadosamente diseñado, surgen costes por la necesidad de controlar y adaptar la ejecución. Comprenden la vigilancia del comportamiento de la contraparte, el ajuste a contingencias no previstas (renegociación, adaptación contractual) y el enforcement o resolución de disputas y se generan una vez que la relación contractual está en marcha.

[Recordemos esta clasificación puesto que nos será muy útil al presentar los diferentes enfoques de la Teoría Contractualista]

Según COASE, toda empresa se enfrenta a un dilema a la hora de decidir si realizar una transacción en el mercado (descentralizada) o si producirla dentro de la empresa de manera interna (centralizada): \textit{to buy vs to make}. Si partimos, por ejemplo, de la existencia de costes de transacción en el mercado, sería más eficiente para una empresa contratar un detenninado servicio o prestación en el mercado, hasta el punto en el que los costes de transacción derivados por realizar una transacción adicional en el mercado igualan a los costes marginales de realizar dicha transacción dentro de la empresa.\footnote{%
    En lenguaje más llano, una empresa debe decidir si contratar un servicio de limpieza provisto por otra empresa o crear un departamento de limpieza dentro de su propia estructura, compuesto por trabajadores en régimen laboral. En principio, esa labor la realizará internamente hasta que el coste marginal de contratar internamente a un limpiador, más igual el coste marginal de contratar dicho servicio en el mercado. A partir de esa necesidad concreta de limpieza (la variable cantidad en este mercado), la empresa subcontrata el servicio externamente a otra empresa.

Otro ejemplo podría verse en una empresa que necesite servicios de transporte urbano y optase entre contratar a un chófer profesional permanente o pagar viajes de taxis a sus empleados. Puede pensarse que para los directivos el coste de ir en taxi cada vez es más elevado, por lo que a la empresa le compensa tener un chofer; mientras que para los empleados el coste de ir en taxi es menor, por lo que a la empresa realiza en el mercado un contrato spot de taxi cada vez que precisa de sus servicios.
} A partir de dicho nivel de provisión del servicio determinado, será más eficiente intemalizar el proceso dentro de la empresa. Se puede partir de la situación inversa en caso de suponer la existencia de costes de transacción dentro de las empresas, utilizando el mismo razonamiento marginal. De hecho, este último ejercicio nos daría la respuesta de Coase a la pregunta del tamaño óptimo de una empresa.

Por tanto, según el autor, las empresas crecen hasta un tamaño determinado en el que el coste marginal de hacer una transacción en la propia empresa supera al de contratarla en el mercado (spot). A partir de ese tamaño, no sería óptimo realizar una transacción centralizada dentro de la empresa:\footnote{%
    Como vemos, las teorías de la empresa mantuvieron el marco de análisis de la eficiencia de la teoría neoclásica (entre otros, el razonamiento marginalista); en cambio, sustituyeron los supuestos irreales del modelo por otros más reales, oponiéndose en cierta medida al modelo de competencia perfecta. Ello implica que todas las teorías a partir de este momento van a asumir la existencia de infonnación imperfecta (incompleta y asimétrica), como ya comenzase a hacer KNIGHT al asumir riesgo (una modalidad de información incompleta). Como sabemos, en caso de información imperfecta, no se garantiza la optimalidad paretiana, siendo el equilibrio de mercado subóptimo.
} las imperfecciones del mercado son la razón de ser de la empresa. Es decir, las teorías de la empresa consideran que ésta es una respuesta eficiente a la información imperfecta, donde la naturaleza del contrato se justifica en la minimización de los costes de transacción que surgen al utilizar factores de producción en el proceso productivo (HOLSTROMM y TIROLE, 1989). El contrato en sí mismo es una disposición de límites dentro de los cuales debe operar el proveedor de bienes o servicios a una empresa (COASE, 1937). 

\paragraph{Teorías contractuales: post-COASE}
Las teorías contractuales de la empresa (post-COASE) se pueden clasificar en dos grupos: en uno de ellos las partes son capaces de realizar contratos completos, es decir, éstos son capaces de regular todas las posibles contingencias futuras en el momento presente. Tales teorías son:
\begin{la}
    \item La empresa como nexo de contratos: ALCHIAN-DEMSETZ (1972); JENSEN-MECKLING (1976);
    \item La teoría del principal-agente o teoría de la agencia: HOLMSTRON-MILGROM (1994).
\end{la}

En el otro grupo de teorías contractuales, las partes sólo pueden realizar contratos incompletos, de modo que, en el momento de negociación del contrato, las partes no pueden prever todas las posibles contingencias de su relación contractual, bien porque es dificil de negociar o porque es muy costoso hacerlo. Ello podría llevar a que, por ejemplo: (i) las partes no firmasen un contrato completo, el cual podría haberles beneficiarles mutuamente; o, (ii) que el resultado futuro de la acción presente de una de las partes podría variar en función de su poder de negociación, lo cual podría no haber estado previsto en el contrato original. En dicha circunstancia, los contratos que se hubiesen firmado se renegociarían una y otra vez. En este contexto de contratos incompletos, es clave que la redacción de los mismos tenga en cuenta los incentivos que pueden introducirse para disciplinar el comportamiento de los agentes ante posibles eventualidades. Estas teorías son:
\begin{la}
    \item La teoría de los costes de transacción: WILLIAMSON (1975, 1985);
    \item La teoría de los derechos de propiedad: GROSSMAN-HART (1986); HART-MOORE (1990).
\end{la}

\subsubsection{Teorías contractuales (I): contratos incompletos}
Las dos vertientes que ahora presentaremos centran su justificación en la reducción de los costes de transacción que se originan ex ante a la perfección del contrato. A pesar de ello, utilizarán recursos diferentes para explicar el origen de éstos y valorar las implicaciones que esto tiene sobre la formación de éstas, y en particular, sobre el tamaño óptimo que deberían tener. 

\paragraph{La Teoría de los costes de transacción (WILLIAMSON)}
Sobre la base de COASE, WILLIAMSON va a considerar que las transacciones económicas pueden llevarse a cabo o bien dentro de la empresa o bien a través del mercado.\footnote{%
    La teoría de los costes de transacción suele identificarse como la teoría de la empresa de WILLIAMSON (1975), siendo ésta la más ligada a las ideas iniciales de COASE (1937). Las ideas de éste han impulsado una escuela económica que se conoce como transaction cost economics. A su vez, estas ideas se enmarcan en la conocida como "nueva economía institucional", la cual, con un sólido fundamento microeconómico, considera la relevancia que tiene el marco organizativo e institucional en las transacciones económicas.
} Por su parte, la empresa es una estructura organizativa que presenta costes de transacción internos (de governance), y cuyas decisiones se guían por un criterio de autoridad o jerarquía. Por el contrario, la transacción en el mercado va a presentar costes de transacción externos y va a guiar sus decisiones por el mecanismo de precios. La pregunta esencial será:

¿Cuál es el tipo y tamaño de organización que minimiza los costes de trasacción?

Par responder, WILLIAMSON va a añadir dos nuevos componentes a la teoría que explican la aparición de costes de transacción:
\begin{lnum}
    \item \textbf{Racionalidad limitada}. Cualquier individuo, a pesar de ser maximizador de su utilidad, se ve sujeto a una serie de restricciones de capacidad para procesar toda la infonnacíón necesaria para decidir racionalmente;\footnote{%
        Basándose en el concepto originario de HERBET SIMON (1957), uno de los padres de la economía del comportamiento (behavioural economics), hace referencia al hecho de que la capacidad humana es limitada para prever y resolver problemas complejos.
    } y,
    \item \textbf{Oportunismo}. Las partes contratantes siguen ante todo su propio interés lo que podría implicar que se comportasen de forma no cooperativa, eludiendo sus compromisos contractuales, haciendo trampas, engañando u ocultando información. El resultado de la transacción económica sería ineficiente desde un punto de vista paretiano. 
\end{lnum}
	
Los costes de transacción se van a originar al intentar mitigar estos dos problemas y su valor va a depender de la eficiencia con la que la empresa pueda realizar transacciones dentro de ella (sin acudir al mercado). Tres características de la transacción van a impactar positivamente en dicho nivel de eficiencia de la empresa, provocando que, si dicha transacción se realizase en el mercado, ésta exigiría la elaboración de contratos muy complejos, los cuales diílcilmente van a poder contemplar todas las posibles eventualidades (esto es, serán incompletos):\footnote{%
    Ello crearía incentivos para que los agentes se comportasen de manera oportunista, ya que su comportamiento no estaría regulado en el contrato.
}
\begin{lnum}
    \item \textbf{Frecuencia}. La frecuencia con que se realicen las transacciones, siendo los costes de transacción externos (esto es, los de transaccionar en el mercado) crecientes con la frecuencia. Al fin y al cabo, si una transacción es muy frecuente para la actividad de una empresa, suele ser mejor internalizarla (por ejemplo, para beneficiarse de economías de escala o de una reducción de los costes de aprendizaje);
    \item Incertidumbre. Incrementa los costes de transacción en el mercado ya que la información incompleta hace dificil o demasiado costosa la elaboración de contratos completos que puedan prever cualquier posible circunstancia. Al aumentar el coste de concluir contratos completos en el mercado, sería más eficiente que la empresa internalizase la actividad en cuestión sometida a incertidumbre; y,
    \item \textbf{Especificidad de los activos implicados}. Un activo es específico cuando requiere una alta inversión para su adquisición y, además, no es susceptible de usos alternativos (o no lo es sin que ello suponga un coste elevado). Por ejemplo, a una empresa de ingeniería civil en una región montañosa le resultaría más eficiente la compra de una tuneladora (y, por lo tanto, que se convierta ésta en parte de la empresa) que el alquiler de los servicios de una empresa de tuneladoras cada vez que necesitase realizar un túnel al construir una carretera.
\end{lnum}

Por supuesto, existen mecanismos de coordinación híbridos entre las dos estructuras extremas (contratos spot en el mercado y pertenencia a la empresa), tales como los contratos de suministro de largo plazo, las redes empresariales, las joint venture o las alianzas estratégicas.  Las fusiones y adquisiciones verticales se originan por el deseo de optimizar las transacciones. En casos intermedios, las empresas pueden optar por alianzas estratégicas o joint-ventures (i.e alianza entre un laboratorio biotecnológico y una farmacéutica). La externalización (outsourcing) de actividades (contabilidad, marketing) puede justificarse porque son actividades no esenciales: frecuentes, pero con bajo grado de especificidad.

La Teoría de los Costes de Transacción también ha sito utilizada para analizar la estructura organizativa más efectiva. Tipos de estructura:
\begin{lnum}
    \item La estructura en U: La actividades de las empresas se subdividen en áreas (marketing, contabilidad, etc) gestionadas por un manager. Muy útil para empresas con un solo producto, ya que permite que el personal se especialice en esa área; y,
    \item La estructura en M: La empresa se estructura en divisiones (por ejemplo, según mercados geográficos o por productos). Es útil para grandes empresas, donde cada división funciona con autonomía y la oficina central se encarga de la estrategia a largo plazo y el control. Una variante es la estructura en H o los holdings, donde una empresa matriz posee otras subsidiarias y que se puede aplicar a empresas multinacionales.
\end{lnum}

\paragraph{Teoría de los Derechos de Propiedad}
Propuesto por GROSSMAN y HART  (1985), este enfoque parte de que la empresa es un conjunto de activos sujetos a una propiedad común,\footnote{%
    Según esta definición, si dos activos distintos tienen el mismo propietario, entonces tendremos una única empresa integrada; si tienen diferentes propietarios, entonces estaríamos ante dos empresas que se relacionan en el mercado. Los derechos que se intercambien pueden incluir el derecho a usar, modificar, transferir o extraer un ingreso del activo en cuestión.
} que pueden ser utilizados por más personas al mismo tiempo, algunas de las cuales poseen derechos de propiedad (empresarios), mientras que otras no (empleados); por lo que examinará los contratos y, por tanto, la empresa desde el punto de vista de los derechos de propiedad de dichos activos.

La teoría acepta los supuestos de WILLIAMSON sobre el comportamiento oportunista de los agentes, pero no limita exclusivamente el carácter incompleto de los contratos a la existencia de racionalidad limitada: aunque las partes puedan prever en un contrato presente las posibles consecuencias futuras, dichas estipulaciones no podrán ser lo suficientemente detalladas como para que, por ejemplo, un juez que interprete una disputa basada en dicho contrato no tuviese ninguna duda a la hora de verificar el cumplimiento del mismo.\footnote{%
    HART, SHLEIFER y VISHNY (1997) aplican esta teoría debate sobre la propiedad pública o privada y En este sentido utilizan los derechos residuales de control para explicar la mayor violencia que existe en las cárceles estadounidenses privadas en comparación con las públicas. En Estados Unidos existen cárceles públicas (propiedad del Gobierno) y cárceles privadas. Con las que el Gobierno firma una serie de contratos que contienen una gran cantidad de provisiones. No obstante, ante la dificultad de fijar todas las condiciones de cómo deben funcionar las cárceles privadas, en los contratos no se hace referencia al grado de formación que deben tener todos los guardias de seguridad. Así, hay un incentivo a que el propietario privado de la prisión decidiese contratar guardias con menor formación para reducir los costes. En dicho caso, si el Gobierno quiere aumentar el nivel de seguridad de dicha prisión, se debería renegociar el contrato. No obstante, según los autores, el propietario siempre tendrá incentivos para desarrollar medidas de ahorro en costes que ponga en peligro la Seguridad Pública. Por ello, si el Gobierno quisiese evitar definitivamente este problema, podría adquirir la presión y prohibir a sugerencia la aplicación de medidas que socaven la seguridad.
}

Por tanto, aunque justifica la existencia de la empresa en base a la incompletitud de los contratos, la Empresa mitiga dichos costes de transacción en tanto que se organiza en torno a relaciones de autoridad, pudiendo imponer una solución a los vacíos contractuales: la propiedad de los activos es una fuente de poder que puede utilizarse cuando los contratos son incompletos (prácticamente siempre).\footnote{%
    Por tanto, a diferencia de WILLIAMSON, cuya teoría se fundamenta sobre el beneficio (o reducción del coste) que una estructura organizativa supone en un mercado plagado de contratos incompletos (y cómo éste aumenta en función de la frecuencia, la especificidad de los activos y la incertidumbre presente en el contrato), el enfoque de los Derechos de propiedad plantea que las relaciones de autoridad presentes en la Empresa (cuyo origen se encuentra en la propiedad de los activos) son precisamente la razón por la cual los costes de transacción pueden ser menores en una estructura organizativa (i.e. Empresa) que en el mercado.
} Por lo tanto, los derechos de propiedad son importantes en esta rama de teorías contractualistas porque influyen en el comportamiento de los individuos, especialmente en presencia de incertidumbre. 

Puede hablarse de dos tipos de derechos de propiedad:
\begin{lnum}
    \item \textbf{Derechos de control especial}. Estos son derechos concretos descritos en un contrato (por ejemplo, el derecho a recibir un dividendo si se posee la acción de una empresa). A diferencia de los siguientes, estos derechos pueden ser cedidos por el propietario a otro agente al estar perfectamente definidos; y,
    \item \textbf{Derechos residuales de control}. Cuando es imposible listar en un contrato todas las acciones/decisiones sobre un activo, queda un “residuo” de decisiones no especificadas. Estos hacen referencia al resto de posibles derechos sobre un bien que no están recogidos específicamente en el contrato, siempre y cuando no contradigan otra normal contractual, legal o consuetudinaria (HART, 1985). Estos derechos le corresponden siempre al propietario del activo.
\end{lnum}

La organización de la actividad económica a través de una empresa permite que ésta sea propietaria de todos los derechos residuales de control sobre los activos fisicos que integran la empresa, la propiedad de los cuáles es la respuesta eficiente a una situación de incertidumbre en la que todos los contratos que puedan firmarse entre los agentes económicos van a ser incompletos (mercado).\footnote{%
    Para entender esto ilustremoslo con un ejemplo. Una empresa de fabricación de lunas de automóvil tiene su oficina y planta de fabricación en Sabadell. Su actividad requiere de activos inmóviles materiales (máquinaria), materiales para el proceso, operarios de fábrica y una serie de personal administrativo (probablemente) y de inmovilizado que “complementa” la actividad principal. Debemos ahora preguntarnos: ¿Qué es lo más importante para este empresario? Bien, es evidente que, para desempeñar correctamente su actividad, lo mínimo imprescindible será la maquinaria directamente relacionada con el proceso de fabricación y el personal que participa directamente en este. Pues bien, es precisamente sobre este conjunto de elementos (subconjunto de elementos de la Empresa) sobre la que éste deseará tener todos los Derechos residuales (hasta el límite legal) puesto que ni puede firmar un contrato que estipule al más mínimo detalle el comportamiento de cada uno de los trabajadores (ni tampoco todas las contingencias posibles en el uso de la maquinaria) de tal forma que minimice todos los riesgos /incertidumbre al que se enfrenta, ni se puede arriesgar a verse afectado por una información o contingencia que perjudique directamente su actividad; por lo tanto, sobre la propiedad residual de los Derechos de estos activos que este individuo constituirá una Empresa (como institución organizada con relaciones de autoridad). Sobre todo el resto de elementos que integran la Empresa, éste podrá o bien recurrir al mercado para cubrir sus necesidades o bien integrarlas dentro de ésta por motivos de eficiencia.
}

Al mismo tiempo, como la posibilidad de renegociar continuamente los contratos incompletos reduce los incentivos de las partes para realizar inversiones en la relación contractual (al afectar a la rentabilidad esperada hold-up problem) los autores defienden que la empresa crecerá hasta hacerse con todos los derechos residuales de aquellos activos cuya utilización / explotación óptima requiere de un grado muy elevado de información, mientras que delegará el resto de contratos “no tan críticos” en el mercado, ya que la propiedad de los derechos residuales permite verificar el cumplimiento de todos aquellos derechos imperfectamente definidos. Para GROSSMAN y HART: la empresa existe y tiene el tamaño que tiene porque la concentración de derechos residuales de control en determinadas manos permite evitar el \textit{hold-up problem} y proporcionar los incentivos óptimos a la inversión específica.\footnote{%
    \textbf{Definición.-} (\textit{Hold-up problem}) En muchas relaciones productivas que se dan en el maercado, cada parte debe realizar inversiones específicas: gastos, aprendizaje, adaptación de activos, etc., que solo tienen valor dentro de esa relación. Como los contratos no pueden prever todas las contingencias futuras ni todos los estados del mundo (son incompletos), muchos contratos deberán ser renegociados expost para determinar cómo usar los activos y cómo repartir los beneficios. El hold-up problem ocurre cuando, tras invertir, una de las partes puede / quiere aprovecharse de la dependencia de la otra porque el contrato no la protege completamente. Como esta información es conocida ex ante cualquiera de las partes (más aún cualquiera que deba realizar la inversión en activos) será consciente de que podrá ser expropiado de parte de sus beneficios esperados de su inversión ex post por lo que, como es normal, invertirá menos de lo óptimo. Como sabemos que, en útima instancia, quien posee los derechos residuales de los activos puede decidir unilateralmente cómo usarlos y, por tanto, tiene mayor poder de negociación en la renegociación posterior, la parte que posee los derechos residuales sabe que nadie puede bloquearle o apropiarse totalmente del valor de su inversión. Por eso, la parte con más interés en la relación tiene incentivos más fuertes para invertir en activos o relaciones específicas, porque captará una porción mayor del beneficio generado por esas inversiones o reducirá los riesgos que son íntrinsecos a la relación contractual.
}

Una de las aplicaciones más sobresalientes de la teoría de los derechos de propiedad es el enfoque ownership-location-internalization (OLI).\footnote{%
    El paradigma ecléctico (DUNNING, 1980) es una teoría económica que también se conoce como el modelo OLI, por las siglas en inglés de “Ownership” (propiedad), “Locational” (localización) e “Internalization” (internalización). Esta teoría como modelo de la inversión extranjera directa (IED), afirma que además de los diversos factores que infieren en la decisión de que una compañía realice IED, propone a las ventajas de la ubicación que explican la naturaleza y destino de IED. Con la expresión ventajas de la ubicación, se refiere a las ventajas de aprovechar los recursos o activos propios de determinado lugar en el extranjero que a una empresa le resultan valiosos por combinar con sus propias ventajas.
Véase: \href{https://static1.squarespace.com/static/562636cfe4b043d43a7492bf/t/59a48341f9a61e1e333c503f/1503953729748/note_oli.pdf}{\textit{The Ownership-Location-Internalization Framework}. RUHL (2017)}
}

\subsubsection{Teorías contractuales (II): contratos completos}
Desde otro enfoque diametralmente opuesto tenemos las teorías contractuales basadas sobre la completitud de los contratos. Diferentes enfoques forman esta rama aunque todos ellos comparten que ala finalidad o razón de ser de la Empresa es reducir los costes de transacción ex post que se puedan originar en las interacciones realizadas bajo el único amparo del funcinamiento del mercado en su conjunto.

Éste nos permitirá, no solo los diferentes costes de transacción ex post, así como las causas que lo originan, sino también entrar a valorar el funcionamiento u organización interna de la empresa de manera que se minimicen éstos. Es decir, mientras que en el enfoque ex ante (primigeneo) de la Teoría contractual, los académicos comparten la idea de que la Empresa, una vez formada, es un ente homogeneo con intereses y objetivos cuyo único fin es minimizar los costes de transacción ex ante y, consecuentemente, maximizar los beneficios; el enfoque ex post considera la empresa como un ente heterogéneo (razón que justifica su existencia como unidad de intercambio organizada de manera independiente) donde existen intereses y objetivos diversos que solo pueden converger (o alinearse débilmente) si dichos intercambios se realizan a través de un canal común: la Empresa.

\paragraph{Teorías Gerencialistas}
El génesis de esta rama de investigación se remonta a la investigación pionera que realizaron los autores,  BERLE y MEANS, en “The Modern Corporation and Private Property” (1932) - nótese que es anterior a la publicación de COASE -, sobre la estructura de propiedad y control en la gran empresa moderna de Estados Unidos.\footnote{%
    Estos autores observaron que en las grandes corporaciones de capital disperso, los propietarios (accionistas) habían perdido el control efectivo de la gestión, que se había desplazado hacia los directivos o managers: lo llamaron la separación entre propiedad y control (separation of ownership and control). De ese hallazgo nace el enfoque gerencialista de la empresa, que sostiene que la empresa ya no actúa como una “máquina de maximización del beneficio de los accionistas” ya que los directivos profesionales toman decisiones autónomas, persiguiendo objetivos propios (crecimiento, estabilidad, poder, reputación), no necesariamente la maximización del beneficio a corto plazo. Por tanto, la empresa moderna se convierte en una institución social y económica con intereses propios, no un mero instrumento de los dueños del capital.
} Estos autores rompen con la visión clásica de la empresa como mero nexo de capital (neoclásica, factores productivos y aprovechamiento de procesos) y abren paso a verla como una organización compleja, con división de funciones y objetivos diferenciados, dando lugar al poder gerencial. Sobre la base de estos autores surgirá lo que hoy se conoce como Teoría Gerencial de la Empresa, en la que destacamos los trabajos de:\footnote{%
    BAUMOL (1959): “Business Behavior, Value and Growth”; WILLIAMSON (1964): “The Economics of Discretionary Behavior”; MARRIS (1964): “The Economic Theory of Managerial Capitalism”; GALBRAIGHT (1952-67): “American Capitalism”, “The New Industrial State”
}
\begin{la}
    \item BAUMOL (1959). Modelo de maximización de ventas. Los directivos buscan maximizar las ventas en vez de los beneficios, dado que los sueldos, la imagen de la empresa, los análisis financieros y estratégico dependen de la evolución de las ventas;
    \item WILLIAMSON (1964). Modelo de preferencia por el gasto. Los Agentes emprenden políticas que maximicen su propia utilidad que son variables ligadas al nivel de gasto de la empresa (sueldo, el poder, prestigio, seguridad o la notoriedad profesional) en lugar de los beneficios de los Principales;
    \item MARRIS (1964). Modelo de maximización del crecimiento. En este modelo, el salario y status de los directores depende del tamaño de los departamentos. Los directores buscan maximizar el crecimiento de la demanda (mediante la diversificación de productos) y el crecimiento de los fondos (para financiar la expansión) a costa de los beneficios; y,
    \item GALBRAIGHT (1952-67). Modelo institucional del poder corporativo. La empresa moderna se convierte en una organización planificada, reduciendo la incertidumbre del mercado mediante, cuyo objetivo no es el beneficio a c/p sino la estabilidad del sistema, el crecimiento sostenido y el mantenimiento del poder organizativo. 
\end{la}

Pero, ¿qué podemos decir sobre todos estos modelos? En todos los casos existen unos costes de transacción ex post que únicamente pueden ser mitigados mediante la alineación de objetivos, es decir, mediante la introducción ex ante de cláusula al contrato (de ahí la completitud de los contratos dentro del marco de la empresa) o incentivos que condicionen la conducta que genera este coste y sometiéndola a la voluntad de quien ostentan la propiedad.

\paragraph{La empresa como un nexo de contratos}
Como apuntan JENSEN y MECKLING, la delimitación precisa de los derechos individuales determina cómo se reparten los costes y las remuneraciones entre los miembros de la empresa, por lo que la empresa puede considerarse un conjunto de contratos: las relaciones contractuales constituyen la esencia de la empresa. 

La teoría no distingue entre transacciones a nivel de mercado y de empresa, porque asume que el mercado y la empresa sólo difieren en la naturaleza de los contratos. Es decir, la única diferencia es que, mientras que dentro de la empresa solo se hacen contratos a largo plazo, en el mercado son más comunes los plazos más cortos (tendencia al spot).

Como veníamos diciendo, esta teoría hace hincapié en los costes de transacción ex post (costes de supervisión, medición), que surgen mediante estipulaciones ex ante (las cláusulas contractuales), en un contexto en el que suponemos que las partes son capaces de firmar contratos completos. JENSEN y MECKLING (1976) afirmarán que las relaciones contractuales conflictivas – razón de ser de la empresa - no solo afectan a los trabajadores, sino también a los proveedores, los consumidores, los acreedores, etc.\footnote{%
    ALCHIAN y DEMSETZ (1972) afirmarían que la causa de la existencia de la empresa es la asimetría de infomación y el comportamiento oportunista de los empleados. En caso de que la información fuera simétrica, los empleados no podrían comportarse de manera oportunista ya que el empleador sería capaz de observar en todo momento que la acción del trabajador se desvía de las estipulaciones del contrato y podría rescindir inmediatamente dicha provisión de servicios.
} En todos estos contratos existe el problema de los costes de control o, más en general, de agencia, con información simétrica, no habría razón para que existiese la empresa, al no ser necesaria una estructura de supervisión.  Según su teoría, el objetivo de la creación de la empresa es minimizar dichos costes, que pueden clasificarse en: 
\begin{lnum}
    \item \textbf{Costes de supervisión del principal}. Costes de control del comportamiento del agente;
    \item \textbf{Costes de compromiso de agencia}. Costes destinados a asegurar al principal que el agente no le causará daños, o en otro caso el agente compensará los daños;
    \item \textbf{Costes residuales}. Pérdidas asumidas por el principal, resultantes de la diferencia entre las decisiones del agente y del principal - ya que el agente no toma las mismas decisiones que habría tomado el principal.
\end{lnum}

El estudio de los costes ex post y, en consecuencia, su modelización o análisis ha experimentado diferentes cambios a lo largo de la segunda mitada del siglo XX, como hemos visto al inicio (Teoría gerencialista o mediante la Ineficiencia X de Lienstein, por ejemplo). Sin embargo, el mejor exponente formal es la Teoría de la Agencia (HOLMSTRÖM, 1979; MILGROM), que plantea identificar, en toda relación contractual la existencia de dos agentes:
\begin{la}
    \item El propietario (“Principal”). Con una función de utilidad definida por sus preferencias; y,
    \item El directivo (“Agente”). Con otra función de utilidad definida por sus propias preferencias, normalmente yuxtapuestas.\footnote{%
        La función de utilidad del directivo puede depender de elementos muy variopintos, algunos coincidentes con los objetivos del Principal y otros no. Por ejemplo, su retribución, sus bonus, la influencia que le otorga dirigir la empresa, el prestigio y tamaño de ésta, etc.
    }
\end{la}

\paragraph{Teoría de la Agencia}
Este enfoque surge a principios de los años 70’s, de la mano de SPENCE y ROSS. No obstante, la versión moderna la debemos a HÖLSTROM.\footnote{%
    Nobel de Economía (2016) junto a HART: “Por sus contribuciones a la Teoría de Contratos”
}

Según HOLMSTRÖM (1979), aunque los resultados de la actividad de los directivos podrían no conocerse con total certeza, las partes son capaces de elaborar un contrato óptimo en el momento inicial que seguiría siendo óptimo durante toda la ejecución del mismo, por lo que no tendría que estar sujeto a renegociaciones constantes. Un contrato óptimo se define como aquél que maximiza los beneficios del principal, realiza un reparto adecuado de los riesgos de la actividad y alinea los incentivos o funciones objetivo de las partes. Los elementos clave de la aportación son entonces:
\begin{lnum}
    \item La empresa se define como un conjunto de recursos (los Agentes) que deben ser incentivados por el Principal a través de un contrato, para alinear los objetivos divergentes de las partes;
    \item El conflicto de intereses se manifiesta como incertidumbre en el contrato.\footnote{%
        La existencia de incertidumbre  o riesgo en el contrato puede manifestarse mediante dos vías: la selección adversa o el riesgo moral. La diferencia principal radica en el momento en el que se estará asumiendo éste por parte del Principal: en el caso de la selección adversa enfrentamos un problema ex ante que condicionará la utilidad esperada del contrato; por el contrario, en el caso del riesgo moral el riesgo se manifestará ex post por lo que lo que se intentará condicionar será la realización de una conducta que sea beneficiosa para el Principal [\authorfont{Ver Tema 3.A.13}]. Por ejemplo, los directivos saben más de la empresa y de su funcionamiento que los propietarios (ventaja informacional)
    } Por ello, las partes tienen que hacer un contrato en el que el sistema salarial ofrecido por el Principal se base en métricas observables de las acciones del Agente que puedan ser controladas;
    \item El Principal deberá diseñar un sistema de incentivos adecuado para conseguir que el Agente persiga los objetivos que él desea; y,
    \item El contrato óptimo será aquel que realice una ponderación óptima entre los objetivos y los riesgos: unos incentivos insuficientes conducirán al Agente a perseguir sus propios intereses, pero unos incentivos sobredimensionados pueden inducir al Agente a tomar excesivos riesgos.\footnote{%
        Un ejemplo muy claro del problema de agencia lo vimos en España con la proliferación del crédito barato antes de la crisis de 2008: los directores de las oficinas de los bancos y cajas españoles eran evaluados en función de los activos que eran capaces de generar, sin que el riesgo asociado a dichos activos fuese lo suficientemente ponderado. Así, los incentivos de los empleados de banca les conducían a conceder el máximo crédito posible, aun cuando fuese de dudosa calidad. Con el estallido de la crisis se demostraría que aquella conducta no estaba alineada con la maximización de beneficios de las entidades.
    }
\end{lnum}

A modo de ejemplo, supongamos un directivo $D$ (agente) que trabaja para el propietario de la empresa $P$ (principal), para el cual debe realizar una acción que equivale a un nivel de esfuerzo $E$, la cual le genera al directivo un coste $C=v(E)$ y al propietario un beneficio $\Pi=b(E)$. Ambas funciones, $b$ y $e$, son crecientes con el nivel de esfuerzo $E$. De ahí el conflicto de interés: $D$ prefiere un a bajo y $P$ prefiere un a alto. 

Sea $w$ la remuneración que recibe $D$ por parte de $P$ cuando $\Pi \ge 0$. Suponemos que tanto $D$ como $P$ no están sometidos a una restricción presupuestaria. Así mismo, suponemos que $P$ es neutral al riesgo y $D$ es averso al riesgo (un supuesto realista en un contexto en el que los propietarios capitalistas de una empresa tienen un mayor acceso a la diversificación de inversiones). En este contexto, la utilidad esperada del propietario es:

$$U_{P} = b(E) - w$$

Mientras que la utilidad esperada del directivo es:

$U_{D} = u(w) - v(E) - \frac{1}{2}\cdot r\cdot  Var(w)$

Donde $r > 0$, mide el grado de aversión al riesgo. Si $b$ es estrictamente cóncava y $v$ es estrictamente convexa, existe un único $E^*$ que maximiza el bienestar conjunto de ambos agentes (al igualar ambas funciones de utilidad, maximiza la distancia $b(E^*)-v(E^*)$). Dicho resultado es un óptimo de Pareto y permite que el directivo no asuma riesgos ($Var(w) = 0$). En este contexto, donde suponemos información simétrica, la remuneración $w$ que condujese a dicho óptimo sería aquella que dependiese de $E$. Debido a que el esfuerzo $E$ no es observable para el propietario (o, en caso de serlo, no es perfectamente definible o medible), la remuneración $w$ no puede depender de $E$. Por ello, la remuneración al asalariado va a depender del beneficio $\Pi$, el cual suponemos que es perfectamente observable al constituir un proxy del esfuerzo realizado por el Agente.

Como muestra el modelo original de HOLMSTRÖM (1979), la inclusión de una parte del salario variable en función de los beneficios alcanzados o la participación de los directivos en el capital de la empresa (por ejemplo, a través de stock options) o pueden ser contratos de dirección óptimos que solucionen los problemas de riesgo moral, ya que, al verse beneficiada la dirección de la buena marcha del negocio, se consiguen alinear sus incentivos con los de la propiedad.

Independientemente de la existencia de un contrato óptimo, se puede enunciar una serie de mecanismos de mercado que también permiten alinear los objetivos de Principal y Agente:
\begin{lnum}
    \item Uno de ellos es el mercado de capitales, que, en principio, evalúa objetivamente la actividad de las empresas participantes (tanto las que coticen en los mercados de renta variable como las que emitan bonos en los mercados de renta fija), por lo que proporciona información a los propietarios para juzgar el rendimiento de los directivos;
    \item Otro mecanismo de control es la propia existencia de un mercado de directivos, en el que éstos compiten por las posiciones, teniendo pues un incentivo para adquirir una buena reputación y poder aspirar a un mejor puesto directivo;\footnote{%
        Además, el mercado de capitales, a través de las OPAs, proporciona un mecanismo eficaz para disciplinar el comportamiento de los directivos. Como ilustra MANNE (1965), en caso de desavenencias entre el directivo y una parte de los propietarios, éstos últimos podrían decidir lanzar una OPA para tomar el control total de la empresa y poder decidir así despedir a los directivos actuales (ello constituiría una OPA hostil, esto es, una OPA que va en contra de las preferencias de los directivos).
    } y,
    \item Por su parte, el gobierno corporativo (corporate governance), entendido como un conjunto de buenas prácticas en las relaciones entre la empresa y su Consejo de Administración, los accionistas y el resto de stakeholders, de manera que se puedan alinear al máximo los intereses de todos los participantes.\footnote{%
        En terminología de Teoría de Juegos, ello implicaría que el Juego del Principal-Agente deja de ser un Juego estratégico (en forma normal) para convertirse en un Juego repetido, done factores como la reputación son clave para iteracionesfuturas.
    } 
\end{lnum}

Pese al interés de la teoría de la agencia, algunos desarrollos posteriores han señalado que es probable que esta teoría exagere los problemas de gestión ineficiente por parte de los directivos, sobre todo porque sus planteamientos son estáticos. En efecto, en un contexto dinámico, la competencia reduce sensiblemente los límites de discrecionalidad de los directivos, por los siguientes motivos:
\begin{lnum}
    \item \textbf{Comparación}: los resultados de la gestión de los directivos pueden ser comparados con los de otras empresas o con los de otros departamentos dentro de la propia empresa;
    \item \textbf{Absorciones}. La no maximización de beneficios reduce el valor de la empresa, lo cual aumenta las posibilidades de que otras organizaciones la adquieran (función disciplinar de las adquisiciones), lo que conduciría a un reemplazamiento del equipo directivo. Su deseo de permanencia en el puesto desincentivará, pues, a los directivos a apartarse en exceso de la maximización de beneficios.
\end{lnum}

No obstante, la separación entre propiedad y dirección tiene mucha más incidencia en las empresas públicas, ya que en éstas las salvedades anteriores no operan. De entre los desarrollos de la teoría de la agencia en este contexto de empresa pública destaca el modelo de NISKANEN, en el que el burócrata maximiza su utilidad, que depende positivamente del presupuesto, dando lugar a un sobredimensionamiento de la empresa pública.\footnote{%
    Todas estas ineficiencias fueron identificadas como fallos del sector público, y a lo largo de los años 80 constituyeron uno de los elementos legitimadores de los procesos de privatización y liberalización que se han vivido desde entonces. Se argumentaba que la privatización de las empresas públicas aumentaría su eficiencia al someterlas a la disciplina del mercado y que la liberalización posterior reduciría las barreras de entrada y de salida generando una competencia potencial que impediría comportamientos estratégicos.
}

\imagenfit{Esquema_NISKANEN_EmpPublica.png}{Esquema de NISKANEN relativo al funcionamiento de la Empresa Pública}{Apuntes Víctor Gutiérrez. Tema 3.A.15}




\clearpage
\section{Teoría sobre el comportamietno de la Empresa en el Mercado}
\puntosclave{
    El enfoque de la Escuela deHarvard va a conllevar que las autoridades de defensa de la competencia vayan a sospechar, de serie, de cualquier mercado en el que haya un reducido número de empresa. Por lo tanto, se van a perseguir una serie de conductas y de situaciones que serán consideradas como anticompetitivas.
    
	La escuela de Chicago va a partir del supuesto de que los mercados se comportan de manera competitiva, excepto cuando existe una intervención pública. Por lo tanto, no realizan la crítica que hacían los autores de Harvard a las situaciones de monopolio o a los comportamientos estratégicos, y tendrán una visión menos partidaria de la intervención pública en los mercados.
    
	La teoría de los mercados contestables de Baumol ha permitido superar la rivalidad entre escuelas al poner el foco en el mantenimiento de una competencia efectiva, independientemente del número de empresas que hay en el mercado, la cual se conseguirá siempre que no haya barreras de entrada o de salida en el mercado. La intervención pública debe estar orientada a la eliminación de dichas barreras.
    
	La forma en la que compiten las plataformas en los mercados digitales ha originado que estos presenten una serie de fuentes de ineficiencia. Tanto las elecciones de contabilidad de Paul como la experiencia de las autoridades de competencia han permitido formular unas recomendaciones regulatorias que se adaptan a las particularidades de estos mercados.
}
	
Hasta ahora hemos estudiado diferentes enfoques que pretenden dar respuesta a la razón de ser de la empresa y, en función de la justificación proporcionada, dar respuesta a los problemas que se suscitan (o pueden suscitarse) en su interior, los conflictos de interés que hay en su seno, los objetivos qué persigue y el tamaño óptimo que debe conseguir para perseguir dicho fin. Sin embargo, sabemos que la Empresa no es un ente aislado  y cuya razón de ser y fin exceden con creces los límites de su propia estructura, proyectándose sobre lo que será ahora nuestro principal objeto de estudio: el mercado.

Si la razón de ser de la Empresa la podemos encontrar en las ineficiencias presentes en el Mercado, debemos estudiar como el desarrollo de esta estructura a la vez alternativa y complementaria influye sobre éste ya que, al margen de ser objeto de muchísima relavancia para el Sector Público, constituye el objeto de análisis finalista de la Empresa. Por tanto, intentaremos dar respuesta a la siguiente pregunta respecto al Mercado. ¿Qué es, cómo se forma y qué consecuencias tiene cada tipo sobre la eficiencia (entendida como grado de competencia efectiva)?

\subsection{Antecedentes históricos}
En 1933, CHAMBERLIN (y ROBINSON, 1933) publicó un reconocido artículo en el que resaltaba la importancia de estudiar la estructura de los mercados en su conjunto, marcando por tanto un antes y un después en la Teoría Económica pues constituye el punto crítico a partir del cual se bifurca la Teoría Microeconómica y el estudio de los mercados, dando el pistoletazo de salida a lo que hoy se conoce como la Teoría de la Organización Industrial. 

En este apartado se hará una referencia a los principales debates que han informado la evolución de esta rama y a sus implicaciones de política económica. En relación con esto último, las políticas más relevantes serán la política de defensa de la competencia y la regulación económica.

\subsubsection{La Escuela de Harvard}
Los economistas neoclásicos, al asumir el comportamiento precio-aceptante de las empresas, no centraron su atención en la elaboración de una Teoría del Mercado, limitando su análisis al estudio de la producción y de los costes. 
La Escuela de Harvard, también conocida como escuela estructuralista, tuvo su apogeo entre los años treinta y setenta del siglo XX, con cabezas como CHAMBERLIN, ROBINSON, MASON o BAIN.
Durante los años previos, finales del siglo XIX y primer cuarto del siglo XX, la economía de Estados Unidos había pasado de estar formada por pequeñas empresas a empresas de gran tamaño. Dicho aumento de la concentración había llevado a un aumento de los comportamientos colusivos entre empresas, por lo que la Sherman Act (1890) como la Clayton Act (1914), fueron la respuesta legal a dichos problemas.\footnote{%
    Esta es la época en la que surgen grandes empresas como la Standard Oil de Rockefeller, US Steel, General Motors, Ford, American Tobacco, General Electric o AT \&T; y de los grandes magnates estadounidenses, como J.K Rockefeller, M. Stanley, Carneig, etc. 

En particular, la existencia de competencia imperfecta se hizo evidente tras la Gran Depresión cuando, ante una caída generalizada de la demanda, los precios se mantuvieron bastante estables.
}
Sin embargo, su aplicación no estuvo exenta de problemas, en especial en lo respectivo a la prohibición de monopolizar el mercado. Éstos fueron perfectamente visibles tras la absolución de U.S. Steel de dichas acusaciones en 1920, cuando el Tribunal estimó que el mero tamaño de la empresa (un 70\% de la capacidad productiva de acero) no implicaba una conducta ilegal de intento de monopolización. Dicho enfoque, que precisa del establecimiento de unas relaciones de causalidad (la llamada rule of reason) supuso una motivación intelectual para los economistas del momento.

\paragraph{Paradigma ECR: Estructura-Conducta-Resultado}
MASON propuso en 1939 la recopilación de información cuantitativa sobre las estructuras de mercado en Estados Unidos (precios, beneficios, cuotas de mercado) para poder identificar relaciones de causalidad entre dichas variables y ser de utilidad en los procesos anti-trust. Basándose en sus descubrimientos empíricos, esta Escuela va a plantear el paradigma ECR, Estructura-Conducta-Resultado, según el cual se establece una relación directa de causalidad entre la estructura competitiva de un mercado, la conducta de las empresas en el mismo y los resultados que se obtengan por las empresas.\footnote{%
    Algunas de las empresas referidas anteriormente, protagonizaron famosos casos antitrust en los tribunales estadounidenses, cuyas sentencias se inspiraron en estos principios estructuralistas. El precedente típico de la Escuela lo constituye la separación (break-up) de la Standard Oil en varias de las seven sisters (ExxonMobil, Chevron y otras 32 empresas), dictada por el Tribunal Supremo en 1911 por considerarla un monopolio en el mercado del petróleo. No obstante, la sentencia que se considera paradigmática de la aplicación del enfoque ECR es la condena de ALCOA, productor de aluminio, en 1945 por monopolización del mercado, donde se utilizó como un argumento relevante el mero tamaño de la empresa. Años más tarde, en 1984, AT\&T, el monopolio regulado de telecomunicaciones, también sería obligado a dividirse en pequeñas partes, en base a argumentos estructuralistas.
} En otras palabras, la concentración empresarial conduciría a la colusión. Veamos los tres componentes por separado:
\begin{la}
    \item \textbf{Estructura}. Este elemento hace referencia al tipo de mercado subyacente (competencia perfecta, competencia imperfecta) y estos autores supondrán que es exógena, dependiendo de: (i) el número y la distribución del tamaño de los vendedores y compradores (de manera que cuanto menor sea el número de empresas en el mercado, mayor es la probabilidad de que aumente el poder de mercado); (ii) las condiciones de entrada; (iii) el grado de diferenciación del producto (a mayor diferenciación, menor sustituibilidad entre productos y mayor capacidad de monopolización); (iv) la tecnología y los costes de producción;
    \item \textbf{Conducta}. La estructura del mercado obliga a las empresas bien a tener un comportamiento competitivo (esto es, precio-aceptante) o bien a comportarse de manera estratégica (coludir en precios con otras empresas, repartirse el mercado, realizar publicidad, erigir barreras de entrada, etc); y,
    \item \textbf{Resultado}. En un mercado competitivo en equilibrio, la demanda igualará a la oferta para aquel volumen de cantidad en el que el $P = CMg$. Por el contrario, en una situación de competencia imperfecta, las empresas pueden obtener beneficios económicos puros - por encima de los costes de producción – que constituyen precisamente el incentivo a mantener su poder de mercado. El resultado para la economía es, por lo tanto, una pérdida de eficiencia al desperdiciarse recursos. 
\end{la}

Dada la relación directa de causalidad que los autores establecen, la política de defensa de la competencia debe basarse en la sospecha de que toda estructura con poder de mercado es producto de comportamientos estratégicos o de barreras naturales a la empresa. Lo esencial, por lo tanto, será cuantificar la estructura del mercado para determinar si ésta permite ejercer poder de mercado, el cual, al generar una asignación ineficiente de recursos, sería dañino per se y, por ello, ilegal. 

En consecuencia, esta Escuela da importancia al cálculo de elasticidades de demanda, de estructuras de costes y, especialmente, a los indicadores de concentración del mercado que permiten conocer la estructura de los mercados, donde destacan: (i) ratio de concentración; y, (ii) índice de HERFINDAHL-HISRCHMANN.

\paragraph{Críticas}
A partir de los años sesenta, los economistas comenzarían a identificar una serie de causalidades inversas que pondrían en tela de juicio los argumentos estructuralistas. Así, la conducta de las empresas podría ser capaz de afectar a la estructura del mercado (por ejemplo, las fusiones de empresas tienden a aumentar la concentración del mercado o la inversión en I+D tiende a afectar a la tecnología y a la estructura de costes óptima de un mercado, pudiendo hacer que un monopolio natural fuese una configuración más eficiente). 

Adicionalmente, el propio resultado que se observe en un mercado puede afectar a la estructura del mismo (un mercado que registre beneficios extraordinarios va a señalizar su atractivo y, por lo tanto, podrá volverse más competitivo a medida que entran nuevas empresas). Tal y como destacó CLARK, la Escuela de Harvard peca de una visión estática de los mercados, mientras que el fenómeno de la competencia debe analizarse desde un punto de vista dinárnico (la acción inicial de la empresa, la respuesta de la competencia, las diferentes estrategias que éstas pueden llevar a cabo).

La Escuela asume que la diferenciación de productos es negativa al aumentar el poder de mercado de las empresas. No obstante, una diferenciación que responde a las preferencias de los consumidores puede aumentar su utilidad.

\subsubsection{Escuela de Chicago}
A partir de los años 60’s, un grupo de autores entre los que destacan STIGLER, DEMSETZ, PELTZMAN o SULLIVAN desarrollaron desde la Universidad de Chicago una serie de postulados que, partiendo del modelo microeconómico neoclásico, van a rechazar la mayoría de las tesis y las recomendaciones políticas de la Escuela de Harvard. Recordemos que, en aquellos años, y de la mano de COASE y del juez POSNER, nacía en la misma alma mater el enfoque Law \& Economics, con el que hay grandes trasvases teóricos.\footnote{%
    Menos relación hay - por el ámbito de aplicación, más que por el contexto ideológico - con las teorías macroeconómicas del monetarismo de FRIEDMAN y de sus seguidores, conocidos como los "Chicago boys" en Latinoamérica.
}

Este grupo de autores, basándose en un enfoque primordialmente teórico, consideran que a largo plazo el modelo relevante para el análisis de cualquier mercado (y de cualquier economía) es el de competencia perfecta. Por lo tanto, y a diferencia de Harvard, van a considerar que el ejercicio del poder de mercado es algo temporal, a no ser que exista una protección legal que permita que determinados sujetos se aprovechen de él de manera sostenida. Dicha concepción neoclásica del modelo económico predominante en la realidad llevaba a una serie de implicaciones sobre la escasa peligrosidad del poder de mercado: 
\begin{lnum}
    \item Los acuerdos colusivos a los que podría llevar dicho poder no son sostenibles en el tiempo, tanto por los elevados costes que implica la aplicación de la colusión en precios y el control de la conducta de los rivales, como por los incentivos perversos que tienen las empresas participantes en los cárteles a romperlos;
    \item Los monopolios naturales no constituyen un fallo de mercado que justifique la intervención pública porque siempre pueden asignarse a un explotador privado a través de una subasta que introduzca competencia en el mercado a la hora de pujar las condiciones de provisión del servicio, asegurándose de que los precios para el consumidor no reflejan el poder de mercado. Esto es lo que se conoce como la "Demsetz competition" (1968); y,
    \item Los beneficios derivados del poder de mercado no pueden ser permanentes debido a los comportamientos de entrada de nuevas empresas en los mercados rentables.
\end{lnum}

En opinión de la Escuela de Chicago, la concentración empresarial otorga una mayor eficiencia a las economías, dado un mejor aprovechamiento de las economías de escala. En este sentido, tanto la tecnología disponible como la libertad de entrada en los mercados van a detenninar la estructura del mercado, garantizando una conducta óptima (competitiva) y unos resultados apropiados.

En general, la Escuela de Chicago va a tener una visión crítica de la intervención del sector público en los mercados. De hecho como se deriva de sus planteamientos, si un monopolio ineficiente se está manteniendo en el tiempo se deberá a las barreras regulatorias que impiden la libre entrada de empresas puesto que, de lo contrario, un monopolio con libre entrada conduce a una asignación eficiente de los recursos. 

Durante aquellos años, los tribunales estadounidenses se hicieron eco de las ideas de Chicago y se aceptaron conductas que no se habrían permitido bajo el enfoque de Harvard (la fusión de General Dynamics Corporation en 1974; o la monopolización de Kodak o IBM), en base al ahorro de costes y a la entrada potencial de empresas que disciplinarían el mercado. Como puede observarse, las decisiones de los jueces han ido dado vaivenes en función de las enseñanzas de la Teoría de la Organización Industrial.

\paragraph{Críticas}
Es una Escuela eminentemente teórica que no produjo por lo general resultados empíricos ni una metodología clara. La aparición de la teoría de juegos años después permitiría dar una base al entendimiento de los comportamientos estratégicos que estos autores tanto negaron, de ahí el supuesto general de que el modelo de competencia perfecta puede utilizarse para analizar cualquier mercado.

Además, consideran que la intervención pública es la única fuente posible de monopolio, cuando la realidad es más compleja y no desarrollaron unas guías políticas de acción concreta para las politicas antitrust.

\subsubsection{Escuela Austríaca}
Esta Escuela constituye una visión radical de la organización industrial, criticando la fundamentación de la teoría neoclásica y de la Escuela de Chicago (con la que erróneamente se identifica en ocasiones). En general, los austriacos critican la metodología económica basada en la optimización del comportamiento de los agentes que buscan soluciones de equilibrio, ya que dichos enfoques ignoran aspectos clave como la incertidumbre. Trataremos dos tradiciones di ferentes.\footnote{%
    Al igual que la Escuela de Chicago, esta Escuela también tiene una derivada macroeconómica, en este caso centrada en las teorías sobre el dinero y la deflación. que no entran en el ámbito del presente tema
}

\paragraph{Enfoque del mercado y del empresario: MISES, HAYEK y KIRZNER}
Como expuso Hayek en "The Use of K.n owledge in Society'' ( 1945), los agentes económicos toman las decisiones en un contexto de incertidumbre (heredero de las ideas de Knight), de desconocimiento de las condiciones iniciales de oferta y demanda, las cuales logran averiguarse a través de la competencia y del rol de los precios como mecanismo de inforrnación.\footnote{%
    En esta Teoría también puede leerse una crítica al sistema económico planificado propuesto por autores marxistas como O. LANGE
}

La competencia es precisamente el proceso dinámico por medio del cual las oportunidades de ganancias no descubiertas en el mundo pasan a conocerse, ya que se señalizan mediante la aparición de beneficios extraordinarios. Por tanto, la competencia no genera un equilibrio en sí, sino un orden o tendencia al equilibrio (lo que el economista austriaco denominó catalaxia). En el equilibrio no existe ya ni competencia ni incertidumbre, pues todos los empresarios ya han explotado, mediante coordinación o arbitraje, todas las oportunidades de ganancia que había en la economía. Para los austriacos, el empresario es un arbitraj ista de oportu nidades. Lo que es relevante para los economistas es pues el análisis de dicha catalaxia u orden, donde el mercado fij a los precios y los intercambios se convierten en un mecanismo de orden espontáneo, los cuales van a ocurrir sin necesidad de que una autoridad central intervenga en el comportamiento de los agentes, que se ven motivados por su propio interés en beneficiarse. 

Las implicaciones de política económica que se derivan del planteamiento austriaco son: La politica económica debe estar encaminada a crear un marco institucionaJ adecuado, que pem1ita maximizar la libertad individual, minimizando toda aquella intervención gubernamental que ponga un corsé a la actividad económica espontánea. • El monopolio no es un elemento negativo, ya que los beneficios extraordinarios que genera se derivan de una mayor utilidad para la socied��d al satisfacer una necesidad que está siendo demandada por los consumidores. De hecho, la legislación antitrust tendría la indeseable consecuencia de reducir el incentivo empresarial en la búsqueda de oportunidades y, así, estaría retrasando la ·introducción de nuevos productos, o incluso podría conducir a la no producción de alguno de ellos. El británico Littlechild afirmaba que, si la alternativa a un producto caro era que no existiese dicho producto por no haber incentivos para su provisión, el bienestar no habría mejorado. Por lo tanto, aunque la fundamentación es distinta, llegan a conclusiones similares a Chicago respecto de la escasa importancia de la concentración mientras exista libre entrada.

Dado que el poder de mercado pennanente lo crea la regulación (por ejemplo, al fijar precios), la intervención pública favorece a las empresas incumbentes y desincentiva la entrada de competidoras. Por ello, proponen la fijación de un marco institucional que pennita a las empresas innovar y buscar el máximo beneficio. El método de fijación de precios públicos Price cap o IPC-X, que Littlechild propuso para British Telecom al privatizarla, iba en esta línea.

\paragraph{Enfoque de la empresa innovadora: SCHUMPETER} 
SCHUMPETER, en su obra principal \textit{Capitalism, Socialism and Democracy} (1942), donde lleva a cabo una crítica a la economía marxista, señala que el valor social de las distintas estructuras de mercado no puede ser analizado empleando una perspectiva estática de la competencia en la que el énfasis se pone en la detenninación del precio del producto y donde los procesos de decisión son inmutables. Para el autor austriaco lo más relevante es el salto de un equilibrio a otro, en el cual tendrán una gran importancia las cuestiones relacionadas con la innovación. La idea central de Schumpeter es que el poder de mercado, en la medida en que pennite al empresario disponer de recursos (los beneficios extraordinarios derivados de dicho poder), crea los incentivos adecuados para que se produzca el proceso innovador en las economías. El poder de mercado crea las condiciones ex ante (acumulación de beneficios para poder costear la innovación) y ex post (certeza de que podrá seguir disfrutando de los beneficios extraordinarios adicionales que le pennita la innovación realizada) para ser emprendedor, ya que algunas inversiones son tan arriesgadas que solo se emprenderían bajo la perspectiva de importantes beneficios futuros. En este sentido, la clave del pensamiento schumpeteriano es el proceso de destrucción creativa, que permite a la economía evolucionar de un equilibrio a otro al empujar a los empresarios a innovar por temor de quedarse rezagados. En la dinámica de destrucción creativa las empresas que no sean capaces de adaptarse a un nuevo contexto competitivo serán expulsadas por otras que hayan podido innovar y estén mejor adaptadas. Sin ninguna duda, la concepción schumpeteriana de los ciclos económicos tiene conexiones con el pensamiento evolucionista en biología.\footnote{%
    La visión schumpeteriana de la economía también ha sido reinterpretada por los teóricos del crecimiento económico en lo que se conoce como "modelos neo-shumpcterianos". los cuales deberían ser tratados en el Tema de modelos de crecimiento endógeno. Bajo estos modelos subyace el supuesto de que hay un continuo de empresas que van renovándose a medida que se producen innovaciones. Por lo tanto. más allá de la visión cíclica aquí expuesta, recordemos que el progreso técnico es. uno de los tres ingredientes del crecimiento económico. Destacamos a los siguientes autores bajo esta escuela de crecimiento: Howitt, Aghion. Acemoglu. Helpman o Grossman.
}

De este planteamiento se pueden obtener dos implicaciones de política económica: • La importancia de la inversión y el fomento de la investigación y desarrollo, donde nos podemos encontrar distintos tipos de medidas: regulatorias (facilitación del marco legal), gasto público (inversión pública, avales), ingresos públicos (desgravaciones fiscales, patent box) o estructurales (conexión universidad-empresa). • La introducción de un sistema de patentes. Desde un punto de vista económico, la concesión de una patente a una invención implica la creación de un monopolio temporal sobre la misma. Dicha barrera de entrada va a crear los incentivos financieros a los que se hace referencia más arriba. Por lo tanto. desde un punto de vista del bienestar social, las patentes suponen una fuente de ganancia de bienestar (al crear incentivos para innovar, y crecer) y una fuente de pérdida de bienestar (al introducir poder de mercado). Wi lliam Nordhaus ( 1 967) modelizó este trade-ojf para calcular la vida óptima de una patente desde el punto de vista del bienestar y, su estimación empírica, la situó en 20 años.\footnote{%
    Cifra que coincide coin la que los principales sistemas jurídicos de patentes recogen, como el TRIPS.
} 

Kenneth Arrow (1962) propuso un planteamiento sobre la innovación en economía que contradice la visión schumpeteriana de la misma. Mientras que para el austriaco la estructura de mercado monopolista va a incentivar la innovación, para Arrow el monopolio suponía el mantenimiento de una actitud conservadora que no creaba los incentivos para innovar, de hecho, podría incluso retrasar el progreso. Para Arrow, es la situación de mayor competencia la que crea incentivos para innovar, con el fin de mantenerse en el mercado con la máxima eficiencia.

\subsection{La Nueva Teoría de la Organización Industrial}
Desde los años ochenta, se ha producido un cambio en la forma en la que se estudia la economía industrial, otorgándose una mayor importancia a la investigación empírica, al igual que en otras disciplinas económicas (comercio, crecimiento).\footnote{%
    Este enfoque ha venido de la mano de múltiples economistas como BRESHANAM, GEROSKI, LAU, ROCHET, LAFFONT, MASKIN o REY. En lo fundamental, el nuevo paradigma considera prioritario el análisis del comportamiento de los agentes y en cómo éste afecta a la estructura del mercado. Desde este punto de vista, la nueva teoría de la organización industrial se encuentra más cerca del paradigma ECR que del enfoque de la escuela de Chicago.
} Este cambio metodológico se vio facilitado por los desarrollos de: (i) la teoría de juegos (la cual dio un fundamento teórico a los compo1tamientos estratégicos); (ii) la econometría (que permitió extender el análisis de las series temporales y los datos de panel); y, (iii) la mayor disponibilidad de datos. BRESHANAN cita tres razones que justifican el cambio de enfoque en organización industrial:
\begin{lnum}
    \item La inobservabilidad de los costes marginales y, por lo tanto, de los mark-ups aplicados por las empresas. En un contexto de incertidumbre, en el que el poder de mercado no es claramente observable, la regulación debe adaptarse al mismo y crear los incentivos adecuados;
    \item Cada industria tiene unas características particulares que exige que deban ser estudiadas por separado; y,
    \item Como se criticaba a la Escuela de Harvard, se dan relaciones de causalidad biunívocas entre la estructura, la conducta empresarial y los resultado.\footnote{%
        Esto es, existen relaciones de correspondencia o causalidad entre cada combinación par de estos elementos, por tanto no pueden analizarse de manera separada y extraer una conclusión claraque se corresponda a la situación que se está experimentando.
    }
\end{lnum}

\subsubsection{Barreras a la entrada y mercados impugnables}
BAUMOL (1982) dio un giro de180 grados a la TOI, cuestionando la necesidad de regular los mercados desde el estudio del comportamiento interno de éste y porponiendo otra perspectiva que tendrá gran influencia en la regulación de los mercados en competencia imperfecta (y en particular aquellos caracterizados por ser industrias de red). Para BAUMOL, el número de empresas presentes en una industria es irrelevante ya que las condiciones de entrada y de salida de empresas son las que verdaderamente condicionan el resultado competitivo. 

Si en un mercado (independientemente del tipo), los potenciales entrantes pueden entrar a competir sin desventajas respecto a los operadores instalados (incumbentes), dicho mercado será impugnable (atacable, contestable). Entonces, en una situación en la que existen beneficios extraordinarios y, por tanto, incentivos a la entrada, si dicho mercado es impugnable, pueden darse estrategias hit \& run  que el incumbente intentará evitar, tendiendo, de manera natural, a un comportamiento competitivo (P = CMg) sin beneficios extraordinarios.\footnote{%
    Una estrategia hit & run podría ser la venta de productos a un precio inferior al precio de mercado actual con el objetivo de robar cuota de mercado, la empresa incumbente temporalmente, apropiándose parte de los beneficios que se estaban generando una vez estos se agoten la empresa. Pero antes saldría del mercado.
} 

Así pues, si se consigue que los mercados sean impugnables y, por lo tanto, que haya un clima permanente de competencia potencial, se puede llegar a la solución competitiva sin que se cumplan los requisitos restrictivos del modelo de competencia perfecta (p.ej., un alto número de empresas, rendimientos constantes a escala). Dentro de las condiciones para que los mercados sean impugnables podemos destacar: 
\begin{lnum}
    \item No puede haber ventajas absolutas en costes. Igualdad de acceso a la mejor tecnología, los mismos inputs y al mismo precio que las empresas incumbentes;
    \item El mercado presenta rendimientos constantes a escala; 
    \item Los consumidores no tienen preferencias entre distintas empresas más allá de las resultantes de los precios o la diferenciación vertical de calidad; y,
    \item Puede haber costes fijos pero no costes hundidos (irreversibles, esto es, no hay barreras de salida). Tampoco hay costes especiales que deban ser soportados por los entrantes y no por las incumbentes.\footnote{%
        Tampoco hay información asimétrica ni comportamientos estratégicos. Por ejemplo, en caso de información asimétrica, las entrantes podrían no tener visibilidad sobre qué hay detrás de una estructura de precios/beneficios concreta, pudiendo las estrategias hit \& run tener riesgo (un precio bajo de la incumbente podría señalizar una ventaja absoluta en costes que desincentivase la entrada)
    }
\end{lnum}

La mayoría de dichas recetas tienen que ver con la no existencia de barreras de entrada. Según BAIN (1956), una barrera de entrada es un conjunto de tecnologías o de características del producto que permite a las empresas incumbentes obtener un beneficio económico en el largo plazo.\footnote{%
    BAIN identificó tres tipos de barreras: economías de escala, diferenciación de productos y ventajas absolutas en costes. STIGLER (1968) criticó este enfoque, especialmente en lo referente a que las economías de escala constituían una barrera de entrada. STIGLER ofreció una definición alternativa: un coste de producción que debe ser soportado exclusivamente por la empresa entrante y no por la incumbente.
}

\paragraph{Implicaciones en Política Económica}
El sector público debe centrarse en eliminar las barreras de entrada y de salida en vez de en el número de empresas que conforman un mercado. El monopolio, por lo tanto, no implica necesariamente una pérdida de bienestar y debemos centrar el estudio en el mantenimiento de la competencia efectiva, la eliminación de los costes hundidos y la garantía del acceso de las empresas en condiciones de igualdad.\footnote{%
    Término que recoge como objetivo la Ley 15/2007 de defensa de la competencia, en vez de hacer referencia al mantenimiento de la “competencia perfecta”.
} Los sectores de industrias de red son el ejemplo paradigmático de esta nueva regulación anti barreras de entrada.\footnote{%
    Los distintos métodos de regulación y liberalización de estos sectores se estudiaran en el tema correspondiente. Estos sectores se caracterizan por la presencia de una red o cuello de botella donde hay una fuerte inversión inicial y, al operar como monopolio natural, actúa como barrera de entrada (por ejemplo, la construcción de la línea de telefonía fija por parte de una empresa incumbente, donde proveedores de servicios de telefonía que quieran entrar al mercado no pueden competir al no tener acceso a la red). En este caso, la regulación de las obligaciones de libre acceso a la red, bajo una serie de condiciones, permiten mantener la contestabilidad del resto de segmentos del mercado, tanto las que se encuentren aguas arriba corno aguas debajo de la red (continuado con el ejemplo, la obligación de arrendar la infraestructura telefónica a nuevos operadores).
}

Además, como es evidente, es más eficiente controlar que se cumplen las obligaciones de libre entrada a la red, que evaluar la situación competitiva de un mercado en función del número de empresas, por lo que presenta un gran avance para las autoridades de la competencia.

\paragraph{Aportaciones de TIROLE en el marco de la contestabilidad de los mercados}
Si bien ya hemos analizado a lo largo del tema la teoría de los two-sided markets de TIROLE, vamos a centramos a continuación en otras tres contribuciones, relacionadas con la contestabilidad de los mercados y las barreras de entrada, que fueron citadas en el Nobel:\footnote{%
    TIROLE fue galardonado con el Premio Nobel de Economía en 2014: “Por sus contribuciones al análisis del poder de mercado y la regulación”. Junto a LAFFONT, es una de las figuras más representativas del Instituto de Economía Industrial y la Escuela de Economía de Toulouse. Si bien ambas instituciones académicas no pueden calificarse aún como una escuela económica, sí que se pueden encontrar rasgos comunes en su trayectoria investigadora: (i) la importancia de los incentivos a la hora de diseñar contratos; (ii) la idea de que los monopolios no se regulan por sí mismos y, a su vez, que los esquemas de regulación que no consideran la información asimétrica son ineficientes y (iii) la aplicación microeconómica de la psicología y la información asimétrica.
}
\begin{lnum}
    \item \textbf{Finanzas y contestabilidad}. Las finanzas y las inversiones pueden ser un factor relevante en el análisis de barreras de entrada. Según TIROLE, las grandes inversiones realizadas por las empresas de mayor tamaño pueden suponer un instrumento estratégico para ahuyentar a potenciales rivales. Dichas inversiones actuarían como costes hundidos irreversibles, al reducir los costes marginales de las empresas inversoras y permitir amenazar a los potenciales rivales con una política agresiva de precios a la baja en caso de entrada en el mercado (dado el mayor capital fijo instalado);
    \item \textbf{Justificación del poder de mercado de un cuello de botella}. Las obligaciones de acceso que propueso BAUMOL podrían perjudicar a la empresa red, pues ponen en peligro su capacidad de obtener beneficios por el servicio que prestan. Por ello, hay incentivos para que la empresa red favorezca a un determinado operador, aguas arriba o abajo, con quien esté relacionado (por ejemplo, vía integración vertical o acuerdos) o que imponga restricciones al resto de operadores (negar el acceso, fijar precios altos, utilizar una tecnología incompatible). De estas formas, la red podría recuperar parte del poder de mercado perdido por la regulación (HART y TIROLE, 1990). Más en general, al ser la concesión de poder de mercado por parte del sector público una situación dificil de justificar, TIROLE planteó motivos por los que la garantía pública del poder de monopolio puede o no ser merecida:
    \vspace{1em}
    \begin{center}
    \setlength{\tabcolsep}{2pt}
    \begin{tabular}{|p{0.30\linewidth}|p{0.30\linewidth}|p{0.30\linewidth}|}
    \hline
     & \textbf{Poder merecido} & \textbf{Poder no merecido} \\
    \hline
    \textbf{Concesiones administrativas}\footnote{%
        Así, si una empresa obtiene una concesión administrativa para, por ejemplo, explotar una ruta de autobuses interurbana (en situación de monopolio legal), ello debería asignarse a través de un proceso competitivo que permita introducir la disciplina de mercado en la concesión del poder de mercado (por ejemplo, a través de una subasta).
    } 
    & {\scriptsize A través de una subasta competitiva}
    & {\scriptsize Monopolio legal gratuito} \\
    \hline
    \textbf{Propiedad intelecual}\footnote{%
        En el caso de que el poder de mercado se manifieste a través de la concesión de un derecho de patente, la empresa patentadora debería serlo en base a una innovación significativa que avance el estado de la ciencia.
    } 
    & {\scriptsize Gran innovación}
    & {\scriptsize Innovación obvia y no nueva} \\
    \hline
    \textbf{Cuello de botella en una industria red}\footnote{%
        Un razonamiento similar podría aplicarse a aquellas empresas que gestionan la red de una industria vertical (como el transporte de electricidad), las cuales deberían ganar dicho derecho en base a un esfuerzo, y no a la mera tradición histórica.
    } 
    & {\scriptsize Esfuerzo inversor}
    & {\scriptsize Ganado por suerte o por condiciones de demanda o costes} \\
    \hline
    \end{tabular}
    \end{center}
    \vspace{1em}
    \item \textbf{Regulación en un contexto de información asimétrica}. En relación con esta última idea, TIROLE ha realizado numerosas contribuciones en la cuestión de cómo regular los precios que fijan las empresas que operan el monopolio natural de una industria de red. Los enfoques tradicionales tendían a fijar unos precios regulados relacionados con los costes de dichas empresas. No obstante, sus aportaciones de se caracterizan por incluir un contexto de información asimétrica, en el que las empresas reguladas no tienen incentivos para revelar sus verdaderos costes, presionando al alza los precios regulados resultantes. Así, Tirole propondrá influyentes ideas como la elaboración de menús de precios y contratos en función del riesgo que deba soportar la red o la comparación de los precios de una empresa con los de otra similar (yardstick competition).
\end{lnum}

\subsection{Las aportaciones de TIROLE al marco de la NTOI}
La idea central en los planteamientos de TIROLE  es que hay pocas recetas de carácter abstracto que sirvan para todos los mercados e industrias dado que ningún mercado es igual. Entre sus principales aportaciones a la economía industrial podemos destacar:
\begin{lnum}
    \item El debate sobre competencia o regulación es artificial y en realidad se presenta en los términos equivocados. La regulación debe ser una herramienta complementaria a la competencia;\footnote{%
        En palabras de TIROLE:” lo que hemos tratado de hacer es crear una regulación que sea suficientemente ligera para que la innovación se lleve a cabo y se promueva la inversión de los antiguos monopolistas. La mala regulación puede reducir notablemente el crecimiento y crear multitud de problemas. Tenemos que crear reglas que promuevan una mejor regulación. Este es nuestro trabajo como economistas: promover una mejor y más eficiente regulación”.
    }
    \item TIROLE y FUNDENBERG, 1984. Caracterizan bajo que condiciones puede ser óptimo para las empresas sobreinvertir o subinvertir en aspectos como la publicidad o la investigación y desarrollo con el objetivo de detener la entrada de competidores;
    \item Sobre el papel del regulador. El regulador puede hacer uso de la información que obtiene sobre el comportamiento de la empresa para posterior regulación. Sin embargo, la empresa puede distorsionar sus acciones para que el regulador interprete la información de manera equivocada. La interacción repetida entre empresa y regulador termina siendo peor para la sociedad que si solo se interactuase una vez (fundamental cuando se diseñan esquemas regulatorios que duran años);
    \item “Competition in telecommunications” (TIROLE y LAFFONT, 2000). Importantes aportaciones en el marco de la liberalización de mercados como el de las telecomunicaciones durante los años 80 y 90. Tratan los peajes de acceso cuando se regula el acceso al componente no competitivo: 
    \begin{la}
        \item Por un lado el peaje debe ser lo suficientemente alto para permitir al antiguo monopolista cubrir sus costes y proporcionar incentivos para que lleve a cabo inversiones. También debe ser los suficientemente bajo como para que los competidores puedan participar en igualdad de condiciones. Por ende el peaje óptimo se compone del coste de la infraestructura + pérdida del antiguo monopolista + incremento del precio cuando la demanda de los entrantes varía poco con el precio; y,
        \item En caso de separación en partes recíprocas: En principio el precio de terminación (peaje de acceso) debería ser irrelevante, pero no lo es cuando las empresas compiten para capatar clientes. Precios de terminación altos reducen los incentivos a captar clientes por lo que se pueden elevar mucho y en este caso se suele recomendar que el precio de terminación sea fijado por el regulador.
    \end{la}
    \item Mercado de las tarjetas de crédito (ROCHET y TIROLE, 2002). El debate está en si las autoridades de competencia deberían intervenir en la fijación de las comisiones que cobran a sus clientes ya que una regulación estricta de las comisiones a comerciantes puede no ser deseable. Los autores muestran que cuando hay más competencia entre los bancos para captar comerciantes que usuarios de tarjeta, una comisión alta puede deshacer el poder de mercado y recuperar la eficiencia; y,
    \item Modelo de TIROLE-LAFFONT (1986). [Ver Tema 3.A.17]
\end{lnum}

\subsubsection{La metodología de la NTOI}
Esta Escuela sobre todo se va a centrar en conocer la conducta de las empresas ya que, ineficiencia de la estructura del mercado no se encuentra en su estructuración, sino en la conducta que se lleva a cabo dentro de él:

\vspace{1em}
\begin{center}\setlength{\tabcolsep}{2pt}
    \begin{tabular}{|p{0.25\linewidth}|p{0.15\linewidth}|p{0.25\linewidth}|p{0.25\linewidth}|}
        \hline
        \textbf{La competencia...} &  & \multicolumn{2}{c|}{\textit{¿Es deseable?}} \\
        \hline
         &  & \textbf{Sí} & \textbf{No} \\
        \hline
        \multirow{2}{*}{\textit{¿Es factible?}} 
        & \textbf{Sí} & {\scriptsize Competencia} & {\scriptsize Exceso de Entrada} \\
        \cline{2-4}
        & \textbf{No} & {\scriptsize Barreras a la Entrada} & {\scriptsize Monopolio Natural} \\
        \hline
    \end{tabular}
\end{center}
\vspace{1em}
 
Esta conducta va a ser estimada de múltiples formas:
\begin{la}
    \item Comparaciones de modelos no anidados (non-nested models). En vez de mirar indicadores estructurales de concentración (que prejuzgan una estructura competitiva concreta), esta metodología consiste en analizar el comportamiento observado de las empresas sin realizar ningún supuesto previo sobre el modelo microeconómico que subyace en dicho mercado;\footnote{%
        Autores como SELLEN, HARSANY, SPENCE o SCHELLING, basándose en las ideas de Nash y Von Neumann y Morgenstern, realizaron múltiples contribuciones que permitieron que cualquier aspecto de la competencia imperfecta se pudiese analizar bajo el prisma de la Teoría de Juegos, dotando al investigador de una mayor caja de herramientas a la hora de observar los mercados. Por ejemplo, estos autores pusieron de manifiesto cómo el grado de acceso a la información que pueda tener un agente puede influir en la probabilidad de que los agentes se comporten de manera colusiva (cooperativa).
    Véase: (i) \href{https://sci-hub.mk/10.2307/2555893}{\textit{Testing static olipoly models: conduct and cost in the sugar industry 1890-1914}. GENOVESE y MULLING, 1998}; y, (ii) \href{https://pages.stern.nyu.edu/~acollard/Bresnahan1955.pdf}{\textit{Competition and Collusion in the American Automobile Industry: The 1955 Price War} BRESNAHAM, 1987}
    } y/o,
    \item Estimar la variación de precios y de cantidades ante variaciones de las funciones de demanda o de costes. Por ejemplo, el grado en el que las empresas reaccionan a cambios en la demanda puede suponer una evidencia directa de su grado de poder de mercado. Se pueden comparar dichas reacciones de las empresas a los cambios que podrían esperarse bajo un modelo de competencia concreto (colusión, Nash, etc.) para ver si éstas son consistentes con dicha estructura de mercado.\footnote{%
        Véase: (i) \href{https://www.sciencedirect.com/science/article/abs/pii/S0167718788800122}{\textit{Estimating the residual demand curve facing a single firm}. BAKER y BRESNAHAM, 1988}; y, (ii) \href{https://onlinelibrary.wiley.com/doi/abs/10.3982/ECTA13333?.com}{\textit{Understanding the Price Effects of the MillerCoors Joint Venture}. WEINBERG y MILLER, 2017}
    }
\end{la}


\clearpage
\section{\textit{Multi-sided Platforms}: Economía digital}
Hoy en día nos encontramos en el mundo digital con una serie de empresas cuya razón de ser dista de los motivos expuestos tanto por la teoría neoclásica como por las contractualistas. Además, la importancia su creciente importancia en nuestras economías ha dado lugar a que cada vez haya más interés en estudiarlas desde una perspectiva tanto teórica como aplicada (política de competencia y regulación). 

Dentro de este ámbito, y motivado por redefinir el concepto de empresa adecuándolo a las caractrísticas de estos mercados, surge el concepto de plataformas o \textit{multi-sided platforms}.\footnote{%
    La modelización de las plataformas tiene su origen en TIROLE y ROCHET (2003), donde exponen la noción de \textit{two-sided markets}, esto es, una plataforma que pondría en común a dos grupos de agentes.

Si bien la teoría se ha generalizado, extendiéndose al concepto de multi-sided markets, platforms o matchmakers (EVANS y SCHMALENSEE, 2016). En el contexto actual de crecimiento de la economía digital, la mayoría de las Big Techs cumplen con las características de una plataforma, si bien el concepto también podría aplicarse a situaciones no digitales. Puede pensarse en una discoteca, una tarjeta de crédito, la moneda de curso legal de un país o una cadena de televisión. En aras de una mayor sistematicidad, se podría plantear una clasificación de distintos tipos de plataformas que no busca ser completa ni general:
\begin{la}
    \item \textbf{Facilitadoras de intercambios}. Determinadas plataformas tienen como principal objetivo poner en contacto a dos agentes para que realicen un intercambio (Wallapop, Tinder, BlaBlaCar, etc);\footnote{El caso de Amazon es representativo de las vertiginosas dinámicas de los mercados digitales donde, en un primer momento, nos encontraríamos con un facilitador de intercambios de libros y, en la actualidad, con un vendedor de productos (donde ya no se darían con tanta intensidad las características puras de plataforma).
    }
    \item \textbf{Financiadas por publicidad}. Son plataformas en las que un grupo de agentes está interesado en mostrar publicidad a otro grupo de agentes y, como veremos a continuación, pagan por ello (Facebook, Telecinco, 20 Minutos); 
    \item \textbf{Sistemas de transacciones}. Permiten realizar transacciones entre distintos grupos de usuarios (Visa, el euro, un criptoactivo); y,
    \item \textbf{Plataformas de software}. Ponen en común a desarrolladores de programas informáticos o Apps con los potenciales usuarios (el sistema operativo Windows, iOS).
\end{la}
} El rasgo característico de este tipo de empresas es que su función principal es la de poner en común a varios grupos de usuarios con distintas características. 

En un contexto de altos costes de transacción, donde a los distintos usuarios les es difícil o costoso encontrar una contraparte para realizar un intercambio económico, la plataforma les puede ayudar a encontrarlo a cambio, o no, de una contraprestación.\footnote{%
    Según el teorema de Coase, en ausencia de fricciones y costes de transacción, el resultado de la negociación entre un oferente y un demandante será óptimo. No obstante, en presencia de altos costes de búsqueda, no se produce un emparejamiento directo del oferente y del demandante, por lo que surge un nuevo actor, la platafonna intennediaria, la cual ayuda a las partes a reducir los costes de transacción. Por ejemplo, la aplicación The Fork reduce los costes de búsqueda para los comensales y reduce los costes de gestión de las reservas para los restaurantes.
}

\subsection{Elementos diferenciadores: Características empresariales}
La caracterización económica de las plataformas difiere de las empresas neoclásicas en varios puntos. 

\subsubsection{\textit{Network-effects}: Efectos de red o externalidades}
En la mayoría de los casos, el valor añadido que aporta la plataforma es satisfacer el deseo de un grupo A de interactuar con otro grupo B. En la medida en que dicho grupo B es más numeroso, el valor de la plataforma aumenta para el grupo A (por ejemplo, un periódico de tirada nacional en vez de provincial, siendo el grupo B los lectores y el grupo A los anunciantes en el periódico).

No obstante, los efectos de red también podrían ser internos al grupo B (los usuarios están en lnstagram debido a que sus amigos también lo están, independientemente del lado A de anunciantes e instagramers) o incluso recíprocos entre el lado A y el lado B (en Wallapop, ambos grupos se benefician de que haya tanto muchos compradores como muchos vendedores). En este sentido, BELLETLAMME (2021) diferencia entre cross-group y with-in group network effects. Para el autor, la gestión de los efectos de red de los distintos grupos de usuarios es la actividad característica de una plataforma.

\subsubsection{\textit{Switching-costs}: Costes para el usuario}
Elevados costes de cambiar de empresa (switching costs). En algunos casos, una vez que el usuario ha empezado a interactuar con una determinada plataforma, le puede ser muy costoso cambiar de plataforma (en términos de tiempo, costes de búsqueda, aprendizaje, etc). Además, los datos personales ya acumulados por una plataforma juegan un papel importante a la hora de retener a sus clientes (por ejemplo, el hecho de tener listas de reproducción creadas en Spotify).

\subsubsection{\textit{Critical mass}: Economías de escala}
En especial en el caso de las plataformas digitales, el coste marginal de añadir un usuario más a la provisión del servicio es cercano a cero. Además, un gran número de plataformas presentan el conocido como problema de la masa crítica, esto es, para que la plataforma pueda crear valor, es necesario que reúna una base significativa de clientes en todos los grupos relevantes. Por lo tanto, deben recurrir a una serie de estrategias para alcanzar dicha masa crítica. 

Las plataformas financiadas mediante publicidad tienden a realizar una actividad de generación de contenido para los usuarios previa a la monetización de los espacios publicitarios. Otra solución podría ser la estrategia de dividir y vencer (JULLIEN, 2011), según la cual la plataforma subsidiaría a los agentes que pertenezcan al grupo con una mayor elasticidad-precio, para que se sumen más fácilmente a la plataforma y, así, una vez la plataforma haya ganado escala de clientes, poder atraer a otros grupos.\footnote{%
    Por ejemplo, Diners Club, la primera empresa emisora de tarjetas de pago en los años 50's, se vio enfrentada al problema de cómo conseguir que los consumidores adoptasen el pago por tarjeta. Inicialmente comenzó distribuyendo tarjetas para los comensales de una serie de restaurantes en Nueva York para que pagasen sus cenas y así se acostumbrasen a esta nueva forma de pago. Una vez se construyó una base sólida de clientes, el negocio se expandió, convirtiéndose en una tarjeta de pagos aceptada en la mayoría de tiendas.
}

\subsubsection{\textit{Big data}: Algoritmos y datos}
Además de las barreras de entrada ya mencionadas (efectos de red, economías de escala, switching costs) es importan te hacer una referencia al papel de los datos como fuente de ventaja competitiva en los mercados digitales. Como apuntan JONES y TONETTI en un influyente artículo (2020),\footnote{\href{}{\textit{Nonrivalry and the Economics of Data}. JONES y TONETTI (2020)}
} los datos presentan características de no rivalidad, esto es, pueden ser utilizados por distintos agentes económicos sin que su dotación disminuya. En otras palabras, el análisis masivo de datos presenta características de bien público. Por lo tanto, hay incentivos por parte de las empresas a acumular datos y a evitar que otras empresas accedan a estos, generando así barreras de entrada significativas, tanto en mercados digitales como tradicionales. Una empresa con amplios datos puede servir mejor las necesidades de los clientes, mejorar sus procesos internos, extraer un mayor excedente del consumidor, etc. Una empresa entrante en un mercado puede tener grandes dificultades para competir sin estos conocimientos que arrojan los datos (\textit{data is the new oil}). Los autores apuntan a que los rendimientos crecientes a escala que presentan los datos permitirían un mayor bienestar social en caso de que los datos fuesen considerados un cuello competitivo (como en las industrias de red tradicionales) y su acceso fuese permitido a todas las empresas que compiten en un mercado. Como veremos, estas ideas han influenciado en gran medida a las nuevas regulaciones en esta materia.  

\subsection{Elementos diferenciadores: Características de mercado}
La economía digital modifica, por tanto, una parte sustancial de los paradigmas tradicionales que han sido objeto de estudio de la organización industrial, a la vez que se mantienen una serie de características que urden sus raíces en el origen de las industrias de red tradicionales.

\subsubsection{\textit{Gatekeepers}: Barreras de entrada}
Una primera cuestión que habría que abordar a la hora de estudiar los mercados digitales es su estructura de mercado. 

\paragraph{Competition \textbf{in} the market}
En principio, se podría afirmar que las plataformas compiten entre sí en una serie de mercados digitales (\textit{competition \textbf{in} the market}). Por ejemplo, Amazon, Google Shopping y Ebay podrían ser competidores en un mercado de búsqueda de compras en internet. Si tenemos en cuenta el negocio de la publicidad en las plataformas, podría argumentarse que muchas de las plataformas existentes (YouTube, Facebook, Instagram,  etc.) compiten por el tiempo de atención de los usuarios a los que se le puede mostrar publicidad. En este sentido, la competencia online podría conceptualizarse como una especie de oligopolio, con fuertes barreras de entrada y comportamientos estratégicos. 

\paragraph{Competition \textbf{for} the market}
No obstante, un rasgo característico de los mercados digitales es que, tanto las características de las plataformas como los comportamientos estratégicos que siguen, pueden dar lugar a que cada platafonna desarrolle un ecosistema que constituya un mercado propio o mercado \textbf{cautivo}. Esto es lo que se denomina como competition FOR the market, de manera que la plataforma que consiga hacerse con el mercado va a generar unos fuertes efectos de red, economías de escala y switching costs que hará que sea poco probable que pierda su posición de dominio (citando a Abba: winner-takes-it-all-competition). Por ejemplo, Apple podría entenderse como un fabricante de hardware, de un sistema operativo, de aplicaciones propias, de una tienda de aplicaciones en la que se venden otras aplicaciones, además de una serie de dispositivos que presentan múltiples complementariedades entre ellos. En estas circunstancias, la empresa puede concebirse como un monopolio que, además, tiene la capacidad de crear e imponer una serie de reglas que van a afectar tanto a los consumidores como a otras empresas (\textit{platforms as regulators} o \textit{gatekeepers}).\footnote{%
    \textit{Big Tech and the Digital Economy: the Moligopoly Scenario}. N. PETIT, 2020. Discute la estructura de mercado subyacente en los mercados digitales, a caballo entre el monopolio y el oligopolio (\textit{moligopoly}) según su tésis.

Piénsese en las estrictas reglas y condiciones contractuales que Apple establece para el Apple Store (única forma para las aplicaciones más pequeñas de acceder, por ejemplo, al mercado de juegos móviles) o las políticas de usuario que pueda diseñar Twitter (téngase en cuenta el caso de la suspensión de la cuenta de Trump y el rol políticodemocrático que pueden jugar determinadas plataformas).
}

\paragraph{Análisis de mercado: Estructura resultante}
Además, el poder de mercado de estas plataformas puede ser dificil de rivalizar, dadas las características ya descritas de fuertes economías de escala, efectos de red y switching costs (existen escasos incentivos para que la iniciativa empresarial cree un buscador alternativo a Google o un chat a WhatsApp que consiga atraer una masa crítica de usuarios). En relación con la estructura de mercado, tiene interés plantear la cuestión de qué rol puede jugar la introducción de algoritmos sobre las estructuras de mercado.\footnote{%
    Ver \href{}{\textit{Virtual Competition: The promise and perils of the algorithm-driven economy}. EZRACHI y STUCKE (2016)}
} EZRACHI y STUCKE plantean que, en abstracto, los algoritmos podrían llevar a un mercado a tres estructuras diferentes:
\begin{la}
    \item \textbf{Monopolio discriminador de primer grado}. Al permitir el análisis masivo de datos conocer perfectamente las características de un consumidor concreto y fijar un precio ajustado que le extraiga todo su excedente;
    \item \textbf{Oligopolio}. Ya que los algoritmos que controlan en tiempo real los precios de las empresas competidoras pueden favorecer la colusión en precios; y,
    \item \textbf{Competencia perfecta}. Pues los algoritmos empleados en favor de los consumidores podrían mejorar la transparencia en los mercados sobre la formación de precios y permitir que se redujesen hasta los costes marginales.
\end{la}

El enfoque regulatorio que se siga en relación con los algoritmos y la inteligencia artificial podrá determinar el equilibrio resultante.

\subsubsection{\textit{Single homing} vs. \textit{Multi-homing}: Relaciones \textit{ad extra}}
Desde otro punto de vista, podríamos establecer dos modalidades distintas en las que los agentes pueden comportarse cuando interactúan con una plataforma. Si los consumidores solamente contratan habitualmente con una plataforma de entre todas las disponibles, se dice que hacen \textit{single-homing} (por ejemplo, un usuario que únicamente se comunica online vía WhatsApp). No obstante, si el usuario tiende a cambiar de plataforma, se dice que hace \textit{multi-homing} (por ejemplo, un usuario que decide entre Uber, Cabify o Bolt en función del precio). 

En muchas plataformas suele darse la situación de que los vendedores (otras empresas) tienden a hacer multi-homing mientras que los compradores (consumidores) tienden a hacer single-homing.\footnote{%
    Por ejemplo, un consumidor tiende a leer el mismo periódico, mientras que un anunciante puede elegir anunciarse en varios periódicos a la vez.
} Esto crea un cuello de botella competitivo, pues la plataforma debe asegurarse tener muchos single-homers para poder atraer multi-homers. Como muestra ARMSTRONG (2006), cuanto más single-homing tiende a hacer un lado del mercado, menor precio va a soportar (incluso puede tener un precio cero o negativo, en caso de que reciba una prestación gratuita). Mientras que el lado multi-homer va a soportar mayores precios.\footnote{%
    Este razonamiento es similar al que se expuso cuando hablamos de \textit{pricing} en multi-sided markets.
} 

\subsection{Implicaciones de Política Económica}
Desde un punto de vista positivo, se argumenta que las plataformas han introducido un elevado grado de innovación en las economías y una mayor utilidad para los consumidores. En este sentido, podría argumentarse que un consumidor prefiere que exista un único buscador monopolista (Google) que, gracias a ser el único repertorio de referencias indexadas, sea capaz de aportar las búsquedas más ajustadas al consumidor, en contraposición a lo que ocurriría si el mercado de búsqueda online estuviese fragmentado en varias empresas. 

Sin embargo, desde un punto de vista normativo, existe un debate sobre las bondades y riesgos que introducen las plataformas que se ha señalado continuamente desde los años 2000's por parte de autoridades de defensa de la competencia. Algunas de estas conductas son:  
\begin{lnum}
    \item \textbf{Precios predatorios o abusivos}. Una decisión óptima de preciospara una plataforma puede ser fijar precios extraordinariamente bajos, incluso negativos, a una parte del mercado (efectos red); mientras que para otra parte sean extraordinariamente altos;\footnote{%
        De manera más general, la política de precios de una platafonna es más compleja que en una empresa neoclásica ya que la demanda de sus servicios no va a depender exclusivamente del precio del bien y/o de otros bienes sustitutivos y/o complementarios: va a depender del número de miembros que haya en otros grupos, en función del tipo de efectos de red que existan. En cierto sentido, los distintos lados del mercado actúan como \textit{bienes complementarios}
    }
    \item \textbf{Foreclosure}. El foreclosure es un tipo de distorsión vertical de acceso en la que una empresa niega el acceso de otra empresa a una parte del mercado a la que necesita acceder para competir. Dicha conducta se ha detectado tradicionalmente en las industrias de red, cuando una empresa verticalmente integrada negaba el acceso de otras comercializadoras de gas a la red de gaseoductos o lo hacía con unas condiciones abusivas, exportando así su poder de monopolio en el mercado de transporte de gas al mercado de la comercialización, al dificultar la competencia de las otras comercializadoras;\footnote{%
        En el mundo digital, podemos ver como Microsoft, empresa fabricante de hardware, al preinstalar en todos sus ordenadores Windows Media Player (tying), dificultaba la competencia en el segmento software de otros desarrolladores de programas de música. El caso Google Android (2018), en el que se condena a Google por preinstalar en los dispositivos Android las aplicaciones de Google, va en el mismo sentido
    }
    \item \textbf{Estándares propios}. El caso paradigmático de este tipo de práctica es el de Apple. Esta empresa ha sido objeto de investigación en numerosas ocasiones por la adopción de cargadores o entradas de conexión distintas, que impedían la competencia en otros mercados conexos con sus dispositivos hardware;
    \item \textit{\textbf{Self-referencing}}. En caso de que una plataforma opere a la vez como plataforma y vendedor, podría darse un conflicto de intereses.
\end{lnum}

\subsubsection{Marco tradicional: Limitaciones}
La respuesta tradicional a los excesos del poder de mercado ha venido de la mano de la política de defensa de la competencia, tanto fuera como dentro del mercado digital.\footnote{%
    Esto es, la persecución de acuerdos ilegales entre empresas o de abusos de posición dominante
} 

Las multas impuestas por la Dirección de Competencia de la Comisión Europea a las Big Techs (también, aunque en menor medida, la Federal Trade Commission estadounidense) tanto por celebración de acuerdos entre empresas (art. 101 TFUE) como por abuso de posición de dominio (art. 102 TFUE) han sido la reacción más notable en los últimos años a los problemas que presentan estos actores. No obstante, dos cuestiones principales ponen de manifiesto la debilidad de este enfoque como paradigma de intervención pública para hacer frente al poder de mercado digital:
\begin{la}
    \item \textbf{Técnicas tradicionales: perjuicio díficil de cuantificar}. Por un lado, la política de defensa de la competencia se basa en una serie de instrumentos analíticos que no recogen bien las características de los mercados digitales. Ya se ha hablado de la distinta visión de los precios predatorios bajo un enfoque neoclásico y bajo un enfoque de economía de plataformas. No obstante, un criterio fundamental para justificar que una conducta empresarial perjudica a los consumidores se basa en el excedente del consumidor. En mercados digitales, donde a veces lo racional es que el precio del consumidor sea cero, ¿cómo puede calcularse el daño que sufre el consumidor? ¿la situación de monopolio de Google está afectando cuantitativamente al bienestar del consumidor? La cuestión probatoria, basada en argumentos económicos y cuantificables, se vuelve compleja en este contexto; y,
    \item Por otro lado la demografía empresarial de los mercados digitales apenas ha variado a pesar de todas las sanciones impuestas a lo largo de los años, y los problemas fundamentales que se han ido detectando siguen ocurriendo y, en cierta medida, magnificándose. De hecho, el problema de la intervención ex post en estos mercados es que, habida cuenta de lo expuesto: una vez que se consigue una posición dominante, es difícil que mediante una multa vaya a revertirse la estructura de ese mercado a una situación de competencia efectiva o contestabilidad.
\end{la} 

\subsubsection{Nueva Teoría de la Organización Industrial: Economía Digital}
En vista de esta situación, se han planteado una serie de respuestas alternativas a las herramientas clásicas de defensa de la competencia:
\begin{lnum}
    \item \textbf{Fragmentación}. La imposición de remedios estructurales a empresas en mercados que presenten problemas estructurales de competencia. Ello podría implicar que una autoridad de competencia obligase a una empresa a fragmentarse, lo que recuerda a las viejas recetas antitrust aplicadas a los monopolios de red de Estados Unidos a principios del siglo XX. En la UE actual, se propuso en 2019 dar dichas competencias a la Comisión, a través de lo que se llamó la New Competition Tool, tal y como lo hace la autoridad británica, si bien se decidió emplear un enfoque regulatorio;
    \item \textbf{Regulación} [\authorfont{Ver Tema 3.A.20}]. Optar por una intervención ex ante en los mercados a través de instrumentos de regulación adecuados. Dentro de estos mecanismos encontramos el Reglamento General de Protección de Datos de la UE (GDPR, 2018) que, entre muchas otras garantías para asegurar la protección de la privacidad de los consumidores, se posibilita la transferencia de datos de una empresa a otra (reducir los switching costs). Más recientemente, destaca la Digital Markets Act de la UE que, dada la dificultad de perseguir abusos de posición dominante por parte de plataformas, pasa a prohibir directamente muchas de las conductas que llevan sancionando las autoridades de competencia estos últimos años.\footnote{%
        Algunas obligaciones positivas como asegurar la portabilidad de datos de los usuarios de una plataforma a otra, permitir la interoperabilidad de hardware y software, permitir la desinstalación de aplicaciones. Y, también, una serie de prohibiciones como vincular contratación de distintos servicios provistos por la plataforma, la prohibición del \textit{self-referencing} o la prohibición de las cláusulas de exclusividad.
    }
\end{lnum} 




 












\clearpage
\section{Conclusión} 

\subsection{Extensiones y relación con otras partes del Temario}


\subsection{Opinión}



\subsubsection{Idea final (Salida o cierra)}


\clearpage
\pagenumbering{gobble}
\MyBibliography
\MyBibItem{Autor}{Título}{Año}{DOI -- LINK}




%ANEXOS
\clearpage
\anexo{\textbf{Preguntas Test}}
%\input{\TemaCode/questions.tex} 




\end{document}